\subsection{Ejemplo 16: Viga articulada con rótula interna y empotramiento final (Práctica 1 - Subgrupo A)}

\begin{tcolorbox}[colback=blue!3!white, colframe=blue!40!black, arc=3pt, boxrule=0.8pt, left=8pt, right=8pt, top=6pt, bottom=6pt]
    \centering
    {\Large\bfseries\color{blue!60!black} Estructura con rótula interna y empotramiento (diagramas $N, Q, M_f$)}\\[2pt]
    {\normalsize\bfseries Viga de $L = 8\text{ m}$ con apoyo móvil en $A$, rótula en $x=5\text{ m}$ y empotramiento en $B$ ($x=8\text{ m}$)}

    \vspace{4pt}\hrule\vspace{4pt}
    {\small\textbf{Unidades de trabajo:} $L\,[\text{m}] \quad F\,[\text{kN}] \quad M\,[\text{kN}\cdot\text{m}] \quad q\,[\text{kN/m}]$}
\end{tcolorbox}

\paso{1}{Datos y clasificación estática}

Consideramos una viga de $L = 8\text{ m}$ apoyada en el nodo $A$ ($x=0$) mediante un carrito móvil, con una rótula o articulación interna en la sección $x = 5\text{ m}$, y empotrada en el extremo derecho $B$ ($x=8\text{ m}$).
Las cargas aplicadas son:
\begin{itemize}
    \item Carga triangular de $0$ a $20\text{ kN/m}$ en el primer tramo de $3\text{ m}$ ($x=0$ a $x=3\text{ m}$).
    \item Momento concentrado antihorario $M_0 = 5\text{ kN}\cdot\text{m}$ en la sección $x = 5\text{ m}$ (coincidiendo con la zona articulada).
    \item Carga distribuida uniforme $q = 10\text{ kN/m}$ a lo largo del tramo final de $3\text{ m}$ ($x=5\text{ m}$ a $x=8\text{ m}$).
\end{itemize}

\begin{figure}[H]
    \centering
    \adjustbox{max width=0.98\textwidth}{%
        \begin{tikzpicture}[scale=1.1, >=Latex]
            % Viga (azul)
            \draw[line width=2.5pt, blue!80!black] (0,0) -- (8,0);

            % Apoyo A (Móvil) en x=0
            \coordinate (A) at (0,0);
            \draw[thick] (A) -- ++(-0.35,-0.35) -- ++(0.7,0) -- cycle;
            \draw[thick] ($(A)+(-0.2,-0.47)$) circle (0.12);
            \draw[thick] ($(A)+(0.2,-0.47)$) circle (0.12);
            \draw[line width=1pt] ($(A)+(-0.5,-0.6)$) -- ($(A)+(0.5,-0.6)$);
            \node[above left=2pt] at (A) {$A$};

            % Rótula en x=5m
            \filldraw[fill=white, draw=black, line width=1pt] (5,0) circle (3pt) node[above=4pt, font=\small\bfseries] {Rótula ($x=5\text{ m}$)};

            % Empotramiento en B (x=8m)
            \coordinate (B) at (8,0);
            \draw[line width=1.5pt] (8,-0.8) -- (8,1.2);
            \foreach \y in {-0.7,-0.4,-0.1,0.2,0.5,0.8,1.1} {
                    \draw[thick] (8,\y) -- (8.25,\y+0.2);
                }
            \node[above left=2pt] at (B) {$B$};

            % Carga triangular entre x=0 y x=3m (h=1.4)
            \draw[line width=1.1pt, red!80!black] (0,0) -- (3,1.4);
            \foreach \x in {0.6, 1.2, 1.8, 2.4, 3.0} {
                    \pgfmathsetmacro{\h}{\x*(1.4/3.0)}
                    \draw[->, line width=0.8pt, red!80!black] (\x, \h) -- (\x, 0.05);
                }
            \node[above, red!80!black, font=\small\bfseries] at (3, 1.45) {$20\text{ kN/m}$};

            % Momento concentrado M = 5 kN*m en x=5m (antihorario)
            \draw[-{Latex[length=2.5mm, width=1.8mm]}, line width=1.2pt, magenta!80!black] ($(5,0)+(40:0.45)$) arc (40:320:0.45);
            \node[above left, magenta!80!black, font=\small\bfseries] at (4.9,0.5) {$5\text{ kN}\cdot\text{m}$};

            % Carga uniforme q = 10 kN/m entre x=5m y x=8m (h=1.0)
            \draw[line width=1.1pt, red!80!black] (5,1.0) -- (8,1.0);
            \foreach \x in {5.0, 5.6, 6.2, 6.8, 7.4, 8.0} {
                    \draw[->, line width=0.8pt, red!80!black] (\x, 1.0) -- (\x, 0.05);
                }
            \node[above, red!80!black, font=\small\bfseries] at (6.5, 1.05) {$10\text{ kN/m}$};

            % Cotas: 3m, 2m, 3m
            \draw[|<->|, >=Latex] (0,-0.9) -- (3,-0.9) node[midway, fill=white, font=\small] {$3\text{ m}$};
            \draw[|<->|, >=Latex] (3,-0.9) -- (5,-0.9) node[midway, fill=white, font=\small] {$2\text{ m}$};
            \draw[|<->|, >=Latex] (5,-0.9) -- (8,-0.9) node[midway, fill=white, font=\small] {$3\text{ m}$};
        \end{tikzpicture}%
    }
    \caption{Esquema de la viga con rótula interna en $x=5\text{ m}$ y empotramiento en el extremo derecho.}
\end{figure}

\paragraph{Clasificación estática del sistema:}
\begin{itemize}
    \item Grados de libertad ($N=1$ barra con rótula libera 1 momento): $3m = 3(1) = 3$.
    \item Reacciones en los apoyos: $r = 4$ ($V_A$ en el apoyo móvil, $V_B, H_B, M_B$ en el empotramiento).
    \item Articulaciones/liberaciones internas: $h = 1$ (rótula en $x=5\text{ m} \implies M(5) = 0$).
    \item Grado de hiperestaticidad:
          \[
              GH = r - 3 - h = 4 - 3 - 1 = 0 \implies \text{\textbf{Estructura isostática.}}
          \]
\end{itemize}

\paso{2}{Reacciones en los apoyos}
Aprovechando la condición de rótula en $x = 5\text{ m}$ ($M_{\text{rótula}} = 0$ para la porción izquierda $A \to \text{Rótula}$):
1. Carga resultante del tramo triangular: $R_q = \frac{1}{2} \cdot 3 \cdot 20 = 30\text{ kN}$ aplicada a $x = 2\text{ m}$ desde $A$.
2. Tomando momentos respecto a la rótula ($x=5\text{ m}$) desde la izquierda:
\[
    \sum M_{\text{rótula}}^{\text{izq}} = 0 \implies V_A \cdot 5 - 30 \cdot (5 - 2) = 0 \implies 5\,V_A - 90 = 0 \implies \boxed{V_A = 18.00\text{ kN}}
\]
3. Equilibrio global de la viga completa:
\begin{align}
    (\rightarrow +) \quad \sum F_x = 0      & \implies \boxed{H_B = 0}                                                                                     \\
    (\uparrow +) \quad \sum F_y = 0         & \implies V_A + V_B - 30 - (10 \cdot 3) = 0 \implies 18 + V_B - 60 = 0 \implies \boxed{V_B = 42.00\text{ kN}} \\
    (\circlearrowleft +) \quad \sum M_B = 0 & \implies M_B + V_A \cdot 8 - 30 \cdot 6 - 5 - (10 \cdot 3) \cdot 1.5 = 0                                     \\
                                            & \implies M_B + 18 \cdot 8 - 180 - 5 - 45 = 0                                                                 \\
                                            & \implies M_B + 144 - 230 = 0 \implies \boxed{M_B = 86.00\text{ kN}\cdot\text{m}}
\end{align}

\paso{3}{Diagramas de esfuerzos internos ($N, Q, M_f$)}

\begin{figure}[H]
    \centering
    \adjustbox{max width=0.98\textwidth}{%
        \begin{tikzpicture}[scale=1.0, >=Latex]
            % Líneas de referencia para alineación vertical
            \draw[dotted, line width=0.7pt, gray] (0, 0.5) -- (0, -7.5);
            \draw[dotted, line width=0.7pt, gray] (3, 0.5) -- (3, -7.5);
            \draw[dotted, line width=0.7pt, gray] (5, 0.5) -- (5, -7.5);
            \draw[dotted, line width=0.7pt, gray] (8, 0.5) -- (8, -7.5);

            % --- DIAGRAMA N(x) ---
            \node[left=6pt, font=\small\bfseries] at (0,0) {$N$};
            \draw[line width=1.2pt] (0,0) -- (8,0);
            \node[above, font=\small] at (4,0) {$N = 0$};

            % --- DIAGRAMA Q(x) ---
            \draw[line width=0.8pt, gray!60] (0,-2.5) -- (8,-2.5);
            \node[left=6pt, font=\small\bfseries] at (0,-2.5) {$Q$};

            % Tramo 1 (0 a 3m): parabólico de +18 a -12 kN (pasa por 0 en x=2.32m)
            % Tramo 2 (3m a 5m): constante en -12 kN
            % Tramo 3 (5m a 8m): lineal de -12 kN a -42 kN
            \fill[blue!15] (0,-2.5) -- (0,-2.5+0.9) to[out=-20, in=150] (2.32,-2.5) -- cycle;
            \fill[blue!15] (2.32,-2.5) to[out=-30, in=170] (3,-2.5-0.6) -- (5,-2.5-0.6) -- (8,-2.5-1.5) -- (8,-2.5) -- cycle;
            \draw[line width=1.3pt, blue!70!black] (0,-2.5+0.9) to[out=-20, in=150] (3,-2.5-0.6) -- (5,-2.5-0.6) -- (8,-2.5-1.5);

            \node[above left, blue!90!black, font=\scriptsize] at (0,-2.5+0.9) {$18.00$};
            \node[below, blue!90!black, font=\scriptsize] at (3,-2.5-0.6) {$-12.00$};
            \node[below, blue!90!black, font=\scriptsize] at (5,-2.5-0.6) {$-12.00$};
            \node[below right, blue!90!black, font=\scriptsize] at (8,-2.5-1.5) {$-42.00\text{ kN}$};

            % Cortante nulo en x = 2.32 m desde A
            \draw[dotted, line width=0.7pt, red!80!black] (2.32, -2.5) -- (2.32, -7.5);
            \draw[|<->|, >=Latex] (0,-4.0) -- (2.32,-4.0) node[midway, fill=white, font=\tiny] {$2.32\text{ m}$};

            % --- DIAGRAMA M_f(x) ---
            \draw[line width=0.8pt, gray!60] (0,-5.5) -- (8,-5.5);
            \node[left=6pt, font=\small\bfseries] at (0,-5.5) {$M_f$};

            % M(0) = 0
            % M(2.32m) = 27.87 kN*m (máximo positivo)
            % M(3m) = 24.00 kN*m
            % M(5m izq) = 0 (Rótula)
            % Salto por momento antihorario de 5 kN*m a la derecha de la rótula: M(5m der) = -5.00 kN*m
            % M(8m) = -86.00 kN*m en el empotramiento
            \fill[blue!15] (0,-5.5) to[out=-35, in=180] (2.32,-5.5-1.4) to[out=0, in=165] (3,-5.5-1.2) -- (5,-5.5) -- (5,-5.5+0.25) to[out=-35, in=170] (8,-5.5+1.5) -- (8,-5.5) -- cycle;
            \draw[line width=1.3pt, blue!70!black] (0,-5.5) to[out=-35, in=180] (2.32,-5.5-1.4) to[out=0, in=165] (3,-5.5-1.2) -- (5,-5.5) -- (5,-5.5+0.25) to[out=-35, in=170] (8,-5.5+1.5);

            \node[below, blue!90!black, font=\scriptsize\bfseries] at (2.32,-5.5-1.4) {$27.87\text{ kN}\cdot\text{m}$};
            \node[below, blue!90!black, font=\scriptsize] at (3,-5.5-1.2) {$24.00$};
            \node[above, blue!90!black, font=\scriptsize] at (5,-5.5) {$0$ (Rótula)};
            \node[above left, blue!90!black, font=\scriptsize] at (8,-5.5+1.5) {$-86.00\text{ kN}\cdot\text{m}$};

            % Criterios de signo para tramos positivo y negativo
            \draw[-{Stealth[length=2mm, width=1.4mm]}, line width=0.8pt, blue!70!black] ($(1.45,-6.3)+(-90:0.25)$) arc (-90:-270:0.25);
            \node[font=\small\bfseries, blue!90!black] at (1.8,-6.3) {$+$};
            \draw[-{Stealth[length=2mm, width=1.4mm]}, line width=0.8pt, blue!70!black] ($(2.15,-6.3)+(-90:0.25)$) arc (-90:90:0.25);
            \draw[-{Stealth[length=2mm, width=1.4mm]}, line width=0.8pt, blue!70!black] ($(6.45,-6.3)+(90:0.25)$) arc (90:270:0.25);
            \node[font=\small\bfseries, blue!90!black] at (6.8,-6.3) {$-$};
            \draw[-{Stealth[length=2mm, width=1.4mm]}, line width=0.8pt, blue!70!black] ($(7.15,-6.3)+(90:0.25)$) arc (90:-90:0.25);
        \end{tikzpicture}%
    }
    \caption{Diagramas de esfuerzo axil ($N$), esfuerzo cortante ($Q$) y momento flector ($M_f$).}
\end{figure}

\comp{
    \textbf{Comprobación sistemática de los diagramas de esfuerzos:}
    \begin{enumerate}
        \item \textbf{Relaciones diferenciales de equilibrio interno:}
              \begin{itemize}
                  \item Tramo 1 ($x \in [0, 3\text{ m}]$, triangular creciente $q(x) = \frac{20}{3}x$):
                        \[
                            \frac{dQ}{dx} = \frac{d}{dx}\left(18 - \frac{10}{3}x^2\right) = -\frac{20}{3}x = -q(x), \qquad \frac{dM_f}{dx} = Q(x)
                        \]
                        En el punto de corte con el eje $x$ ($x = 2.32\text{ m}$, donde $Q=0$), se anula la pendiente $\frac{dM_f}{dx} = 0$, obteniéndose el máximo relativo $M_{f,\text{máx}} = +27.87\text{ kN}\cdot\text{m}$.
                  \item Tramo 2 ($x \in [3, 5\text{ m}]$, $q = 0$): Cortante constante $Q(x) = -12.00\text{ kN}$ y momento flector lineal decreciente, verificando $\frac{dM_f}{dx} = -12.00 = Q(x)$.
                  \item Tramo 3 ($x \in [5, 8\text{ m}]$, $q = 10\text{ kN/m}$ uniforme): Cortante lineal decreciente de $-12.00\text{ kN}$ a $-42.00\text{ kN}$ con pendiente $\frac{dQ}{dx} = -10 = -q(x)$ y momento flector parabólico con $\frac{dM_f}{dx} = Q(x)$.
              \end{itemize}

        \item \textbf{Condiciones de contorno en los extremos (reacciones) y rótula:}
              \begin{itemize}
                  \item Extremo izquierdo $A$ ($x = 0$): El cortante arranca en $Q(0) = +18.00\text{ kN}$, coincidiendo con la reacción $V_A = 18.00\text{ kN}$ hacia arriba. $M_f(0) = 0$ por ser apoyo simple.
                  \item Rótula interna ($x = 5\text{ m}$): El momento flector se anula por la izquierda de la articulación: $M_f(5^-) = 0$.
                  \item Empotramiento derecho $B$ ($x = 8\text{ m}$): El cortante cierra en $Q(8) = -42.00\text{ kN}$ (equilibrado por $V_B = 42.00\text{ kN}$ ascendente) y el momento flector cierra en $M_f(8) = -86.00\text{ kN}\cdot\text{m}$, coincidiendo en módulo y sentido con el momento de reacción $M_B = 86.00\text{ kN}\cdot\text{m}$.
              \end{itemize}

        \item \textbf{Discontinuidades y saltos por acciones concentradas:}
              \begin{itemize}
                  \item En $x = 5\text{ m}$ actúa un momento concentrado antihorario $M_0 = 5\text{ kN}\cdot\text{m}$. Provoca un salto negativo brusco en el momento flector desde $0$ (en la rótula) hasta $-5.00\text{ kN}\cdot\text{m}$ ($\Delta M_f = -5\text{ kN}\cdot\text{m}$).
                  \item No hay fuerzas verticales concentradas intermedias, por lo que el cortante $Q(x)$ es continuo en todo el dominio.
              \end{itemize}

        \item \textbf{Grado polinómico según la carga aplicada:}
              \begin{itemize}
                  \item Tramo 1 ($q(x)$ lineal): $Q(x)$ es parabólico (grado 2) y $M_f(x)$ es cúbico (grado 3).
                  \item Tramo 2 ($q = 0$): $Q(x)$ es constante (grado 0) y $M_f(x)$ es lineal (grado 1).
                  \item Tramo 3 ($q = \text{cte}$): $Q(x)$ es lineal (grado 1) y $M_f(x)$ es parabólico (grado 2).
              \end{itemize}
    \end{enumerate}
}

\subsection{Ejemplo 17: Viga continua sobre 3 apoyos con rótula interna (Práctica 1 - Subgrupo B)}

\begin{tcolorbox}[colback=blue!3!white, colframe=blue!40!black, arc=3pt, boxrule=0.8pt, left=8pt, right=8pt, top=6pt, bottom=6pt]
    \centering
    {\Large\bfseries\color{blue!60!black} Viga continua articulada sobre 3 apoyos (diagramas $N, Q, M_f$)}\\[2pt]
    {\normalsize\bfseries Viga de $L = 11\text{ m}$ con apoyos en $A, B, D$, rótula en $x=5\text{ m}$ y carga trapezoidal a triangular}

    \vspace{4pt}\hrule\vspace{4pt}
    {\small\textbf{Unidades de trabajo:} $L\,[\text{m}] \quad F\,[\text{kN}] \quad M\,[\text{kN}\cdot\text{m}] \quad q\,[\text{kN/m}]$}
\end{tcolorbox}

\paso{1}{Datos y clasificación estática}

Consideramos una viga de $L = 11\text{ m}$ apoyada sobre tres apoyos verticales: apoyo articulado en $A$ ($x=0$), apoyo móvil intermedio en $B$ ($x=3\text{ m}$) y apoyo móvil final en $D$ ($x=11\text{ m}$). En la sección $x = 5\text{ m}$ se encuentra una rótula o articulación interna.
Las cargas aplicadas son:
\begin{itemize}
    \item Carga triangular que sube desde $0$ en $A$ ($x=0$) hasta $20\text{ kN/m}$ en el apoyo intermedio $B$ ($x=3\text{ m}$).
    \item Carga uniforme $q = 20\text{ kN/m}$ entre $x=3\text{ m}$ y la rótula en $x=5\text{ m}$ ($2\text{ m}$).
    \item Carga triangular descendente desde $20\text{ kN/m}$ en la rótula ($x=5\text{ m}$) hasta $0$ en el apoyo final $D$ ($x=11\text{ m}$, tramo de $6\text{ m}$).
\end{itemize}

\begin{figure}[H]
    \centering
    \adjustbox{max width=0.98\textwidth}{%
        \begin{tikzpicture}[scale=1.0, >=Latex]
            % Viga (azul)
            \draw[line width=2.5pt, blue!80!black] (0,0) -- (11,0);

            % Apoyo A (Fijo) en x=0
            \coordinate (A) at (0,0);
            \draw[thick] (A) -- ++(-0.35,-0.5) -- ++(0.7,0) -- cycle;
            \draw[line width=1pt] (-0.5,-0.5) -- (0.5,-0.5);
            \foreach \x in {-0.35,-0.1,0.15,0.4} \draw (\x,-0.5) -- (\x-0.15,-0.7);
            \node[above left=2pt] at (A) {$A$};

            % Apoyo B (Móvil) en x=3m
            \coordinate (B) at (3,0);
            \draw[thick] (B) -- ++(-0.35,-0.35) -- ++(0.7,0) -- cycle;
            \draw[thick] ($(B)+(-0.2,-0.47)$) circle (0.12);
            \draw[thick] ($(B)+(0.2,-0.47)$) circle (0.12);
            \draw[line width=1pt] ($(B)+(-0.5,-0.6)$) -- ($(B)+(0.5,-0.6)$);
            \node[below left=2pt] at (B) {$B$};

            % Rótula en x=5m
            \filldraw[fill=white, draw=black, line width=1pt] (5,0) circle (3pt) node[above=4pt, font=\small\bfseries] {Rótula ($x=5\text{ m}$)};

            % Apoyo D (Móvil) en x=11m
            \coordinate (D) at (11,0);
            \draw[thick] (D) -- ++(-0.35,-0.35) -- ++(0.7,0) -- cycle;
            \draw[thick] ($(D)+(-0.2,-0.47)$) circle (0.12);
            \draw[thick] ($(D)+(0.2,-0.47)$) circle (0.12);
            \draw[line width=1pt] ($(D)+(-0.5,-0.6)$) -- ($(D)+(0.5,-0.6)$);
            \node[above right=2pt] at (D) {$D$};

            % Carga distribuida (roja):
            % Tramo 1 (0 a 3m): sube de 0 a 1.5
            % Tramo 2 (3m a 5m): constante en 1.5
            % Tramo 3 (5m a 11m): baja de 1.5 a 0
            \draw[line width=1.1pt, red!80!black] (0,0) -- (3,1.5) -- (5,1.5) -- (11,0);

            \foreach \x in {0.6, 1.2, 1.8, 2.4, 3.0} {
                    \pgfmathsetmacro{\h}{\x*(1.5/3.0)}
                    \draw[->, line width=0.8pt, red!80!black] (\x, \h) -- (\x, 0.05);
                }
            \foreach \x in {3.6, 4.2, 4.8} {
                    \draw[->, line width=0.8pt, red!80!black] (\x, 1.5) -- (\x, 0.05);
                }
            \foreach \x in {5.6, 6.5, 7.4, 8.3, 9.2, 10.1} {
                    \pgfmathsetmacro{\h}{(11-\x)*(1.5/6.0)}
                    \draw[->, line width=0.8pt, red!80!black] (\x, \h) -- (\x, 0.05);
                }
            \node[above, red!80!black, font=\small\bfseries] at (4, 1.55) {$20\text{ kN/m}$};

            % Cotas: 3m, 2m, 6m
            \draw[|<->|, >=Latex] (0,-0.9) -- (3,-0.9) node[midway, fill=white, font=\small] {$3\text{ m}$};
            \draw[|<->|, >=Latex] (3,-0.9) -- (5,-0.9) node[midway, fill=white, font=\small] {$2\text{ m}$};
            \draw[|<->|, >=Latex] (5,-0.9) -- (11,-0.9) node[midway, fill=white, font=\small] {$6\text{ m}$};
        \end{tikzpicture}%
    }
    \caption{Esquema de la viga continua con 3 apoyos y rótula en $x=5\text{ m}$.}
\end{figure}

\paragraph{Clasificación estática del sistema:}
\begin{itemize}
    \item Grados de libertad ($N=1$ barra articulada): $3m = 3(1) = 3$.
    \item Reacciones en apoyos: $r = 4$ ($V_A, H_A$ en el apoyo fijo $A$; $V_B$ en $B$; $V_D$ en $D$).
    \item Liberación por articulación interna: $h = 1$ (rótula en $x=5\text{ m} \implies M(5) = 0$).
    \item Grado de hiperestaticidad:
          \[
              GH = r - 3 - h = 4 - 3 - 1 = 0 \implies \text{\textbf{Estructura isostática.}}
          \]
\end{itemize}

\paso{2}{Diagrama de Cuerpo Libre (DCL) y Descomposición Estructural}

Sustituimos las ligaduras por sus correspondientes reacciones de enlace:
\begin{itemize}
    \item Apoyo fijo $A$ ($x = 0$): reacción horizontal $H_A$ y reacción vertical $V_A$.
    \item Apoyo móvil intermedio $B$ ($x = 3\text{ m}$): reacción vertical $V_B$.
    \item Apoyo móvil extremo $D$ ($x = 11\text{ m}$): reacción vertical $V_D$.
    \item Rótula interna en $C$ ($x = 5\text{ m}$): libera el momento flector interno ($M_f(5) = 0$), transmitiendo únicamente un esfuerzo cortante vertical interno de interacción mutua $V_C$.
\end{itemize}

\begin{figure}[H]
    \centering
    \adjustbox{max width=0.98\textwidth}{%
        \begin{tikzpicture}[scale=1.1, >=Latex]
            % Viga aislada
            \draw[line width=2.5pt, blue!80!black] (0,0) -- (11,0);
            \fill (0,0) circle (2.2pt) node[above left=2pt, font=\bfseries] {$A$};
            \fill (3,0) circle (2.2pt) node[below=2pt, font=\bfseries] {$B$};
            \filldraw[fill=white, draw=black, line width=1.2pt] (5,0) circle (3pt) node[above=4pt, font=\bfseries] {$C$ (Rótula)};
            \fill (11,0) circle (2.2pt) node[above right=2pt, font=\bfseries] {$D$};

            % Líneas discontinuas divisorias de tramo
            \draw[dashed, gray!80, line width=0.8pt] (3,0) -- (3,1.5);
            \draw[dashed, gray!80, line width=0.8pt] (5,0) -- (5,1.5);

            % Carga distribuida exterior
            \draw[line width=1.1pt, red!80!black] (0,0) -- (3,1.5) -- (5,1.5) -- (11,0);

            \foreach \x in {0.6, 1.2, 1.8, 2.4, 3.0} {
                    \pgfmathsetmacro{\h}{\x*(1.5/3.0)}
                    \draw[->, line width=0.75pt, red!80!black] (\x, \h) -- (\x, 0.05);
                }
            \foreach \x in {3.6, 4.2, 4.8} {
                    \draw[->, line width=0.75pt, red!80!black] (\x, 1.5) -- (\x, 0.05);
                }
            \foreach \x in {5.6, 6.5, 7.4, 8.3, 9.2, 10.1} {
                    \pgfmathsetmacro{\h}{(11-\x)*(1.5/6.0)}
                    \draw[->, line width=0.75pt, red!80!black] (\x, \h) -- (\x, 0.05);
                }
            \node[above, red!80!black, font=\footnotesize\bfseries] at (4, 1.55) {$20\text{ kN/m}$};

            % Reacciones en apoyos
            \draw[->, >=Latex, line width=1.4pt, green!60!black] (-1.1,0) -- (-0.05,0) node[midway, above=1pt, font=\footnotesize\bfseries] {$H_A$};
            \draw[->, >=Latex, line width=1.4pt, green!60!black] (0,-1.0) -- (0,-0.05) node[midway, left=2pt, font=\footnotesize\bfseries] {$V_A$};
            \draw[->, >=Latex, line width=1.4pt, green!60!black] (3,-1.0) -- (3,-0.05) node[midway, right=2pt, font=\footnotesize\bfseries] {$V_B$};
            \draw[->, >=Latex, line width=1.4pt, green!60!black] (11,-1.0) -- (11,-0.05) node[midway, right=2pt, font=\footnotesize\bfseries] {$V_D$};

            % Cotas
            \draw[|<->|, >=Latex] (0,-1.3) -- (3,-1.3) node[midway, fill=white, font=\footnotesize\bfseries] {$3\text{ m}$};
            \draw[|<->|, >=Latex] (3,-1.3) -- (5,-1.3) node[midway, fill=white, font=\footnotesize\bfseries] {$2\text{ m}$};
            \draw[|<->|, >=Latex] (5,-1.3) -- (11,-1.3) node[midway, fill=white, font=\footnotesize\bfseries] {$6\text{ m}$};
        \end{tikzpicture}%
    }
    \caption{Diagrama de Cuerpo Libre (DCL) sustituyendo las ligaduras por reacciones.}
    \label{fig:ejemplo13_dcl}
\end{figure}

\paso{3}{Cálculo de las Reacciones en los Apoyos}

Descomponemos el cálculo aislando las dos subestructuras independientes conectadas por la articulación interna:

\begin{enumerate}
    \item \textbf{Subestructura derecha isostática independiente ($CD$, $x \in [5\text{ m}, 11\text{ m}]$):}
          Esta chapa se comporta como una viga simplemente apoyada entre la rótula $C$ y el apoyo móvil $D$.
          \begin{itemize}
              \item La carga triangular desciende desde $20\text{ kN/m}$ en $C$ hasta $0$ en $D$ a lo largo de una base de $6\text{ m}$.
              \item Resultante equivalente:
                    \[
                        R_2 = \frac{20\text{ kN/m} \cdot 6\text{ m}}{2} = 60\text{ kN}
                    \]
              \item Centro de gravedad: a un tercio de la base medido desde el extremo de mayor carga ($C$):
                    \[
                        x_{R2} = 5 + \frac{1}{3}(6) = 7\text{ m} \quad (\text{brazo a la rótula: } 2\text{ m}, \quad \text{brazo al apoyo } D: 4\text{ m})
                    \]
          \end{itemize}
          Equilibrio de momentos respecto a la rótula $C$ por la derecha:
          \begin{equation}
              \sum M_C^{\text{der}} = 0 \implies V_D \cdot 6 - 60 \cdot 2 = 0 \implies 6\,V_D = 120 \implies \boxed{V_D = 20.00\text{ kN}}
          \end{equation}
          Equilibrio vertical de la chapa derecha para determinar la interacción transmitida por la rótula:
          \begin{equation}
              \sum F_y = 0 \implies V_{C,\text{der}} + V_D - 60 = 0 \implies V_{C,\text{der}} = 60 - 20 = 40.00\text{ kN}
          \end{equation}
          Por el principio de acción y reacción, la rótula transmite a la chapa izquierda una fuerza descendente hacia abajo de $40\text{ kN}$ en $x = 5\text{ m}$.

    \item \textbf{Subestructura izquierda ($ABC$, $x \in [0, 5\text{ m}]$):}
          Sobre esta porción actúan:
          \begin{itemize}
              \item Carga triangular en $AB$ ($3\text{ m}$): resultante $R_1 = \frac{20 \cdot 3}{2} = 30\text{ kN}$ aplicada a $x = \frac{2}{3}(3) = 2\text{ m}$.
              \item Carga uniforme en $BC$ ($2\text{ m}$): resultante $R_{\text{rec}} = 20 \cdot 2 = 40\text{ kN}$ aplicada a $x = 3 + 1 = 4\text{ m}$.
              \item Fuerza vertical transmitida por la rótula en $C$: $V_C = 40\text{ kN}$ hacia abajo en $x = 5\text{ m}$.
          \end{itemize}
          Equilibrio de momentos respecto a la rótula $C$ por la izquierda:
          \begin{equation}
              \sum M_C^{\text{izq}} = 0 \implies V_A \cdot 5 + V_B \cdot 2 - 30 \cdot (5 - 2) - 40 \cdot (5 - 4) = 0
          \end{equation}
          \begin{equation}
              5\,V_A + 2\,V_B - 90 - 40 = 0 \implies 5\,V_A + 2\,V_B = 130
          \end{equation}

    \item \textbf{Equilibrio global de toda la viga:}
          \begin{align}
              \sum F_x = 0 & \implies \boxed{H_A = 0}                                                                           \\
              \sum F_y = 0 & \implies V_A + V_B + V_D - 30 - 40 - 60 = 0 \implies V_A + V_B + 20 = 130 \implies V_A + V_B = 110
          \end{align}
          Resolviendo el sistema:
          \begin{equation}
              5\,V_A + 2(110 - V_A) = 130 \implies 3\,V_A + 220 = 130 \implies 3\,V_A = -90 \implies \boxed{V_A = -30.00\text{ kN}}
          \end{equation}
          El signo negativo indica que la reacción en $A$ actúa en sentido descendente (apoyo sometido a tiro o tracción vertical):
          \begin{equation}
              V_B = 110 - (-30) \implies \boxed{V_B = 140.00\text{ kN}}
          \end{equation}
\end{enumerate}

\paso{4}{Método de Cortes (Leyes de Esfuerzos Internos)}

La estructura presenta discontinuidades analíticas y apoyos en $x = 3\text{ m}$ y $x = 5\text{ m}$, por lo que dividimos la viga en \textbf{tres tramos de corte}:

\tramobadge{Tramo 1}{$0 \le x \le 3\text{ m}$}

Aislamos la porción izquierda seccionada a distancia $x \in [0, 3\text{ m}]$. Por semejanza geométrica, la carga triangular a distancia $x$ tiene una altura:
\begin{equation}
    \frac{20}{3} = \frac{q(x)}{x} \implies q(x) = \frac{20}{3}x
\end{equation}
La resultante parcial es $R_1(x) = \frac{1}{2} x \cdot \left(\frac{20}{3}x\right) = \frac{10}{3}x^2$, con centro de gravedad a $\frac{x}{3}$ del corte.

\begin{figure}[H]
    \centering
    \adjustbox{max width=0.85\textwidth}{%
        \begin{tikzpicture}[scale=1.4, >=Latex]
            \draw[line width=2.5pt, blue!80!black] (0,0) -- (3.5,0);
            \draw[thick, fill=gray!20] (3.5,-0.45) to[out=60,in=-60] (3.5,0) to[out=120,in=-120] (3.5,0.45);
            \filldraw[black] (3.5,0) circle (2.2pt);

            \draw[line width=1.1pt, orange!90!black] (3.5,0) -- (3.5,1.5);
            \draw[line width=1.1pt, red!80!black] (0,0) -- (3.5,1.5);
            \foreach \xx in {0.7, 1.4, 2.1, 2.8} {
            \pgfmathsetmacro{\hh}{\xx*(1.5/3.5)}
            \draw[-{Latex[length=1.6mm, width=1.1mm]}, line width=0.6pt, red!70!black] (\xx, \hh) -- (\xx, 0.05);
            }
            \node[above right=2pt, red!80!black, font=\footnotesize\bfseries] at (3.5,1.5) {$q(x) = \frac{20}{3}x$};

            \draw[|<->|, >=Latex] (0,-0.75) -- (3.5,-0.75) node[midway, fill=white, font=\footnotesize\bfseries] {$x$};

            % Reacción en A (hacia abajo)
            \draw[->, >=Latex, line width=1.4pt, green!60!black] (0,0) -- (0,-1.0) node[midway, left=2pt, font=\footnotesize\bfseries] {$V_A = 30\text{ kN}$};
            \draw[->, >=Latex, line width=1.4pt, green!60!black] (-1.0,0) -- (-0.05,0) node[midway, above=1pt, font=\footnotesize\bfseries] {$H_A = 0$};

            % Esfuerzos internos en el corte
            \coordinate (corte) at (3.8,0);
            \draw[->, line width=1.3pt, red!80!black] (corte) -- ++(1.1,0) node[right, font=\footnotesize\bfseries] {$N(x)$};
            \draw[->, line width=1.3pt, red!80!black] (corte) -- ++(0,-1.0) node[below, font=\footnotesize\bfseries] {$Q(x)$};
            \draw[-{Latex[length=2.3mm, width=1.6mm]}, line width=1.3pt, red!80!black] ($(corte)+(-40:0.55)$) arc (-40:120:0.55);
            \node[right=2pt, red!80!black, font=\footnotesize\bfseries] at ($(corte)+(40:0.75)$) {$M_f(x)$};
        \end{tikzpicture}%
    }
    \caption{Diagrama de esfuerzos internos en el Tramo 1 ($0 \le x \le 3\text{ m}$).}
    \label{fig:ejemplo13_corte1}
\end{figure}

Ecuaciones de equilibrio:
\begin{itemize}
    \item \textbf{Esfuerzo axil ($N$):}
          \begin{equation}
              \sum F_x = 0 \implies \boxed{N(x) = 0}
          \end{equation}
    \item \textbf{Esfuerzo cortante ($Q$):}
          \begin{equation}
              \sum F_y = 0 \implies -30 - \frac{10}{3}x^2 - Q(x) = 0 \implies \boxed{Q(x) = -30 - \frac{10}{3}x^2\text{ [kN]}}
          \end{equation}
          Valores en extremos:
          \begin{align*}
              Q(0) & = -30.00\text{ kN}                                    \\
              Q(3) & = -30 - \frac{10}{3}(9) = -30 - 30 = -60.00\text{ kN}
          \end{align*}
    \item \textbf{Momento flector ($M_f$):}
          \begin{equation}
              \sum M_{\text{corte}} = 0 \implies M_f(x) - (-30)x + \left(\frac{10}{3}x^2\right)\left(\frac{x}{3}\right) = 0
          \end{equation}
          \begin{equation}
              \boxed{M_f(x) = -30\,x - \frac{10}{9}x^3\text{ [kN}\cdot\text{m]}}
          \end{equation}
          Valores en extremos:
          \begin{align*}
              M_f(0) & = 0                                                                     \\
              M_f(3) & = -30(3) - \frac{10}{9}(27) = -90 - 30 = -120.00\text{ kN}\cdot\text{m}
          \end{align*}
\end{itemize}

\tramobadge{Tramo 2}{$3\text{ m} \le x \le 5\text{ m}$}

En este tramo intermedio entre el apoyo $B$ y la rótula $C$, consideramos la sección $x \in [3\text{ m}, 5\text{ m}]$:
\begin{itemize}
    \item En $x = 0$: $V_A = -30\text{ kN}$ hacia abajo.
    \item Carga triangular total de $AB$: $30\text{ kN}$ en $x = 2\text{ m}$.
    \item Reacción en $B$ ($x = 3\text{ m}$): $V_B = 140\text{ kN}$ hacia arriba.
    \item Carga uniforme constante $q = 20\text{ kN/m}$ en la longitud $(x - 3)$: resultante $20(x - 3)$ en su punto medio.
\end{itemize}

Ecuaciones de equilibrio:
\begin{itemize}
    \item \textbf{Esfuerzo cortante ($Q$):}
          \begin{equation}
              \sum F_y = 0 \implies -30 - 30 + 140 - 20(x - 3) - Q(x) = 0
          \end{equation}
          \begin{equation}
              \boxed{Q(x) = 80 - 20(x - 3) = 140 - 20\,x\text{ [kN]}}
          \end{equation}
          Valores en extremos:
          \begin{align*}
              Q(3^+) & = 140 - 20(3) = +80.00\text{ kN} \quad (\text{salto brusco en } B: \Delta Q = +80 - (-60) = 140\text{ kN}) \\
              Q(5)   & = 140 - 20(5) = +40.00\text{ kN}
          \end{align*}
    \item \textbf{Momento flector ($M_f$):}
          \begin{equation}
              \sum M_{\text{corte}} = 0 \implies M_f(x) - (-30)x + 30(x - 2) - 140(x - 3) + 20(x - 3)\frac{x - 3}{2} = 0
          \end{equation}
          \begin{equation}
              M_f(x) = -30\,x - 30(x - 2) + 140(x - 3) - 10(x - 3)^2
          \end{equation}
          \begin{equation}
              \boxed{M_f(x) = -120 + 80(x - 3) - 10(x - 3)^2\text{ [kN}\cdot\text{m]}}
          \end{equation}
          Valores en extremos:
          \begin{align*}
              M_f(3) & = -120.00\text{ kN}\cdot\text{m}                             \\
              M_f(5) & = -120 + 80(2) - 10(4) = -120 + 160 - 40 = 0\text{ (Rótula)}
          \end{align*}
\end{itemize}

\tramobadge{Tramo 3}{$5\text{ m} \le x \le 11\text{ m}$}

Para maximizar la simplicidad analítica, evaluamos la porción derecha mirando desde el extremo libre con apoyo $D$ ($x = 11\text{ m}$) hacia la izquierda, introduciendo la coordenada auxiliar medida desde $D$:
\begin{equation}
    u = 11 - x \quad \implies \quad 0 \le u \le 6\text{ m}
\end{equation}
A una distancia $u$ desde el extremo $D$, por semejanza triangular con la carga triangular que sube de $0$ a $20\text{ kN/m}$ en $6\text{ m}$:
\begin{equation}
    \frac{20}{6} = \frac{q(u)}{u} \implies q(u) = \frac{10}{3}u
\end{equation}
La carga resultante sobre la porción derecha es:
\begin{equation}
    R(u) = \frac{1}{2} u \cdot \left(\frac{10}{3}u\right) = \frac{5}{3}u^2 \quad \text{con centro de gravedad a } \frac{u}{3} \text{ desde el corte}
\end{equation}

\begin{figure}[H]
    \centering
    \adjustbox{max width=0.85\textwidth}{%
        \begin{tikzpicture}[scale=1.4, >=Latex]
            \draw[line width=2.5pt, blue!80!black] (0,0) -- (4.5,0);
            \draw[thick, fill=gray!20] (0,-0.45) to[out=60,in=-60] (0,0) to[out=120,in=-120] (0,0.45);
            \filldraw[black] (0,0) circle (2.2pt);
            \fill (4.5,0) circle (2.2pt) node[above right=2pt, font=\bfseries] {$D$};

            \draw[line width=1.1pt, orange!90!black] (0,0) -- (0,1.5);
            \draw[line width=1.1pt, red!80!black] (0,1.5) -- (4.5,0);
            \foreach \xx in {0.8, 1.7, 2.6, 3.5} {
            \pgfmathsetmacro{\hh}{(4.5-\xx)*(1.5/4.5)}
            \draw[-{Latex[length=1.6mm, width=1.1mm]}, line width=0.6pt, red!70!black] (\xx, \hh) -- (\xx, 0.05);
            }
            \node[above left=2pt, red!80!black, font=\footnotesize\bfseries] at (0,1.5) {$q(u) = \frac{10}{3}u$};

            \draw[|<->|, >=Latex] (0,-0.75) -- (4.5,-0.75) node[midway, fill=white, font=\footnotesize\bfseries] {$u = 11 - x$};

            % Reacción en D
            \draw[->, >=Latex, line width=1.4pt, green!60!black] (4.5,-1.0) -- (4.5,-0.05) node[midway, right=2pt, font=\footnotesize\bfseries] {$V_D = 20\text{ kN}$};

            % Esfuerzos internos en la cara izquierda del corte
            \coordinate (corte) at (-0.3,0);
            \draw[->, line width=1.3pt, red!80!black] (corte) -- ++(-1.0,0) node[left, font=\footnotesize\bfseries] {$N(x)$};
            \draw[->, line width=1.3pt, red!80!black] (corte) -- ++(0,1.0) node[above, font=\footnotesize\bfseries] {$Q(x)$};
            \draw[-{Latex[length=2.3mm, width=1.6mm]}, line width=1.3pt, red!80!black] ($(corte)+(40:0.55)$) arc (40:-140:0.55);
            \node[left=2pt, red!80!black, font=\footnotesize\bfseries] at ($(corte)+(-140:0.75)$) {$M_f(x)$};
        \end{tikzpicture}%
    }
    \caption{Diagrama de esfuerzos internos en el Tramo 3 evaluado por la derecha ($5\text{ m} \le x \le 11\text{ m}$).}
    \label{fig:ejemplo13_corte3}
\end{figure}

Equilibrio de la porción derecha:
\begin{itemize}
    \item \textbf{Esfuerzo cortante ($Q$):}
          \begin{equation}
              \sum F_y = 0 \implies Q(x) - \frac{5}{3}u^2 + 20 = 0 \implies Q(x) = \frac{5}{3}u^2 - 20 = \boxed{\frac{5}{3}(11 - x)^2 - 20\text{ [kN]}}
          \end{equation}
          Valores en extremos:
          \begin{align*}
              Q(5)  & = \frac{5}{3}(6)^2 - 20 = 60 - 20 = +40.00\text{ kN} \\
              Q(11) & = \frac{5}{3}(0)^2 - 20 = -20.00\text{ kN}
          \end{align*}
          Posición de cortante nulo ($Q = 0$) en el vano derecho:
          \begin{equation}
              \frac{5}{3}u_0^2 - 20 = 0 \implies u_0^2 = 12 \implies u_0 = \sqrt{12} \approx 3.464\text{ m}
          \end{equation}
          \begin{equation}
              x_0 = 11 - u_0 = 11 - 3.464 = \boxed{7.536\text{ m}} \quad (\text{a } 2.536\text{ m a la derecha de la rótula})
          \end{equation}
    \item \textbf{Momento flector ($M_f$):}
          Tomando momentos respecto a la sección de corte (sentido horario positivo):
          \begin{equation}
              \sum M_{\text{corte}} = 0 \implies -M_f(x) + 20\,u - \left(\frac{5}{3}u^2\right)\left(\frac{u}{3}\right) = 0
          \end{equation}
          \begin{equation}
              \boxed{M_f(u) = 20\,u - \frac{5}{9}u^3 = 20(11 - x) - \frac{5}{9}(11 - x)^3\text{ [kN}\cdot\text{m]}}
          \end{equation}
          Valores en extremos:
          \begin{align*}
              M_f(u=6) & = 20(6) - \frac{5}{9}(216) = 120 - 120 = 0\text{ (Rótula en } x = 5\text{ m)} \\
              M_f(u=0) & = 0\text{ (Apoyo final } D\text{ en } x = 11\text{ m)}
          \end{align*}
          Valor máximo positivo en $u_0 = \sqrt{12}$:
          \begin{equation}
              M_{f,\text{máx}} = 20\sqrt{12} - \frac{5}{9}(12\sqrt{12}) = \left(20 - \frac{20}{3}\right)\sqrt{12} = \frac{40}{3}\sqrt{12} \approx \boxed{46.19\text{ kN}\cdot\text{m}}
          \end{equation}
\end{itemize}


\paso{5}{Diagramas de Esfuerzos Internos ($N, Q, M_f$)}

\begin{figure}[H]
    \centering
    \adjustbox{max width=0.98\textwidth}{%
        \begin{tikzpicture}[scale=1.05, >=Latex]
            % Líneas de referencia para alineación vertical
            \draw[dotted, line width=0.7pt, gray] (0, 1.8) -- (0, -7.5);
            \draw[dotted, line width=0.7pt, gray] (3, 1.8) -- (3, -7.5);
            \draw[dotted, line width=0.7pt, gray] (5, 1.8) -- (5, -7.5);
            \draw[dotted, line width=0.7pt, gray] (7.54, 1.8) -- (7.54, -7.5);
            \draw[dotted, line width=0.7pt, gray] (11, 1.8) -- (11, -7.5);

            % --- ESQUEMA DCL / REACCIONES Y CARGAS ---
            \draw[line width=2.4pt, blue!80!black] (0,1.3) -- (11,1.3);

            % Carga distribuida
            \draw[line width=0.9pt, red!80!black] (0,1.3) -- (3,2.2) -- (5,2.2) -- (11,1.3);
            \foreach \x in {0.6, 1.2, 1.8, 2.4, 3.0} {
            \pgfmathsetmacro{\h}{1.3 + \x*(0.9/3.0)}
            \draw[-{Latex[length=1.7mm, width=1.1mm]}, line width=0.7pt, red!80!black] (\x, \h) -- (\x, 1.35);
            }
            \foreach \x in {3.6, 4.2, 4.8} {
            \draw[-{Latex[length=1.7mm, width=1.1mm]}, line width=0.7pt, red!80!black] (\x, 2.2) -- (\x, 1.35);
            }
            \foreach \x in {5.6, 6.5, 7.4, 8.3, 9.2, 10.1} {
            \pgfmathsetmacro{\h}{1.3 + (11-\x)*(0.9/6.0)}
            \draw[-{Latex[length=1.7mm, width=1.1mm]}, line width=0.7pt, red!80!black] (\x, \h) -- (\x, 1.35);
            }
            \node[above, red!80!black, font=\scriptsize\bfseries] at (4, 2.25) {$20\text{ kN/m}$};

            % Reacciones en apoyos
            \draw[->, >=Latex, line width=1.3pt, green!60!black] (0,1.3) -- (0,0.4) node[midway, left=2pt, font=\footnotesize\bfseries] {$V_A = 30\text{ kN}$};
            \draw[->, >=Latex, line width=1.3pt, green!60!black] (3,0.4) -- (3,1.25) node[midway, right=2pt, font=\footnotesize\bfseries] {$V_B = 140\text{ kN}$};
            \draw[->, >=Latex, line width=1.3pt, green!60!black] (11,0.4) -- (11,1.25) node[midway, right=2pt, font=\footnotesize\bfseries] {$V_D = 20\text{ kN}$};
            \draw[->, >=Latex, line width=1.3pt, green!60!black] (-0.9,1.3) -- (-0.05,1.3) node[midway, above=1pt, font=\footnotesize\bfseries] {$H_A = 0$};

            % --- DIAGRAMA N(x) ---
            \node[left=6pt, font=\small\bfseries] at (0,0) {$N$};
            \draw[line width=1.2pt] (0,0) -- (11,0);
            \node[above, font=\small] at (5.5,0) {$N(x) = 0$};

            % --- DIAGRAMA Q(x) ---
            \draw[line width=0.8pt, gray!60] (0,-2.5) -- (11,-2.5);
            \node[left=6pt, font=\small\bfseries] at (0,-2.5) {$Q$};

            % Curvas de cortante
            \fill[blue!15] (0,-2.5) -- (0,-2.5-0.6) to[out=-20, in=170] (3,-2.5-1.3) -- (3,-2.5) -- cycle;
            \fill[blue!15] (3,-2.5) -- (3,-2.5+1.8) -- (5,-2.5+0.9) to[out=-15, in=150] (7.54,-2.5) -- cycle;
            \fill[blue!15] (7.54,-2.5) to[out=-30, in=175] (11,-2.5-0.5) -- (11,-2.5) -- cycle;

            \draw[line width=1.3pt, blue!70!black] (0,-2.5-0.6) to[out=-20, in=170] (3,-2.5-1.3);
            \draw[line width=1.3pt, blue!70!black] (3,-2.5+1.8) -- (5,-2.5+0.9) to[out=-15, in=150] (7.54,-2.5) to[out=-30, in=175] (11,-2.5-0.5);

            \node[below left, blue!90!black, font=\scriptsize\bfseries] at (0,-2.5-0.6) {$-30.00$};
            \node[below left, blue!90!black, font=\scriptsize\bfseries] at (3,-2.5-1.3) {$-60.00$};
            \node[above left, blue!90!black, font=\scriptsize\bfseries] at (3,-2.5+1.8) {$+80.00$};
            \node[above right, blue!90!black, font=\scriptsize\bfseries] at (5,-2.5+0.9) {$+40.00$};
            \node[below right, blue!90!black, font=\scriptsize\bfseries] at (11,-2.5-0.5) {$-20.00\text{ kN}$};

            % Criterios de signo para cortante
            \draw[-{Latex[length=1.6mm, width=1.1mm]}, line width=0.7pt, blue!80!black] (1.1,-2.5-0.35) -- (1.1,-2.5-0.7);
            \node[font=\small\bfseries, blue!80!black] at (1.35,-2.5-0.52) {$-$};
            \draw[-{Latex[length=1.6mm, width=1.1mm]}, line width=0.7pt, blue!80!black] (1.6,-2.5-0.7) -- (1.6,-2.5-0.35);

            \draw[-{Latex[length=1.6mm, width=1.1mm]}, line width=0.7pt, blue!80!black] (3.8,-2.5+0.35) -- (3.8,-2.5+0.7);
            \node[font=\small\bfseries, blue!80!black] at (4.05,-2.5+0.52) {$+$};
            \draw[-{Latex[length=1.6mm, width=1.1mm]}, line width=0.7pt, blue!80!black] (4.3,-2.5+0.7) -- (4.3,-2.5+0.35);

            \draw[|<->|, >=Latex] (5,-4.0) -- (7.54,-4.0) node[midway, fill=white, font=\tiny\bfseries] {$2.54\text{ m}$};

            % --- DIAGRAMA M_f(x) --- (En color naranja según convenio)
            \draw[line width=0.8pt, gray!60] (0,-5.5) -- (11,-5.5);
            \node[left=6pt, font=\small\bfseries] at (0,-5.5) {$M_f$};

            % Zona negativa en [0, 5m] (dibujada hacia arriba por tracción en fibra superior)
            \fill[orange!20] (0,-5.5) to[out=20, in=-170] (3,-5.5+1.8) to[out=-20, in=170] (5,-5.5) -- cycle;
            \draw[line width=1.3pt, orange!85!black] (0,-5.5) to[out=20, in=-170] (3,-5.5+1.8) to[out=-20, in=170] (5,-5.5);

            % Zona positiva en [5, 11m] (dibujada hacia abajo por tracción en fibra inferior)
            \fill[orange!20] (5,-5.5) to[out=-30, in=180] (7.54,-5.5-1.1) to[out=0, in=165] (11,-5.5) -- cycle;
            \draw[line width=1.3pt, orange!85!black] (5,-5.5) to[out=-30, in=180] (7.54,-5.5-1.1) to[out=0, in=165] (11,-5.5);

            \node[above, orange!90!black, font=\scriptsize\bfseries] at (3,-5.5+1.8) {$-120.00\text{ kN}\cdot\text{m}$};
            \node[above, orange!90!black, font=\scriptsize] at (5,-5.5) {$0$ (Rótula)};
            \node[below, orange!90!black, font=\scriptsize\bfseries] at (7.54,-5.5-1.1) {$+46.19\text{ kN}\cdot\text{m}$};

            % Flechas de signo momento negativo en [0, 5m] dentro del área coloreada
            \draw[-{Stealth[length=2mm, width=1.4mm]}, line width=0.8pt, orange!85!black] ($(1.3,-5.5+0.6)+(90:0.25)$) arc (90:270:0.25);
            \node[font=\small\bfseries, orange!90!black] at (1.65,-5.5+0.6) {$-$};
            \draw[-{Stealth[length=2mm, width=1.4mm]}, line width=0.8pt, orange!85!black] ($(2.0,-5.5+0.6)+(90:0.25)$) arc (90:-90:0.25);

            \draw[-{Stealth[length=2mm, width=1.4mm]}, line width=0.8pt, orange!85!black] ($(7.0,-5.5-0.55)+(-90:0.25)$) arc (-90:-270:0.25);
            \node[font=\small\bfseries, orange!90!black] at (7.35,-5.5-0.55) {$+$};
            \draw[-{Stealth[length=2mm, width=1.4mm]}, line width=0.8pt, orange!85!black] ($(7.7,-5.5-0.55)+(-90:0.25)$) arc (-90:90:0.25);
            \draw[-{Stealth[length=2mm, width=1.4mm]}, line width=0.8pt, orange!85!black] ($(7.7,-5.5-0.55)+(40:0.25)$) arc (40:180:0.25);
        \end{tikzpicture}%
    }
    \caption{Diagramas de esfuerzo axil ($N$), esfuerzo cortante ($Q$) y momento flector ($M_f$).}
    \label{fig:ejemplo13_diagramas}
\end{figure}

\comp{
    \textbf{Comprobación sistemática de los diagramas de esfuerzos:}
    \begin{enumerate}
        \item \textbf{Relaciones diferenciales de equilibrio interno:}
              \begin{itemize}
                  \item Tramo 1 ($x \in [0, 3\text{ m}]$, triangular creciente): $\frac{dM_f}{dx} = -30 - \frac{10}{3}x^2 = Q(x)$ y $\frac{dQ}{dx} = -\frac{20}{3}x = -q(x)$.
                  \item Tramo 2 ($x \in [3, 5\text{ m}]$, uniforme $20\text{ kN/m}$): $\frac{dM_f}{dx} = 80 - 20(x - 3) = Q(x)$ y $\frac{dQ}{dx} = -20 = -q(x)$.
                  \item Tramo 3 ($x \in [5, 11\text{ m}]$, triangular decreciente): $\frac{dM_f}{dx} = \frac{5}{3}u^2 - 20 = Q(x)$. En el punto de cortante nulo ($x = 7.54\text{ m}$), $\frac{dM_f}{dx} = 0$, alcanzándose el máximo positivo $M_{f,\text{máx}} = +46.19\text{ kN}\cdot\text{m}$.
              \end{itemize}

        \item \textbf{Condiciones de contorno en los extremos (reacciones) y rótula:}
              \begin{itemize}
                  \item Extremo izquierdo $A$ ($x = 0$): Cortante inicial $Q(0) = -30.00\text{ kN}$, coincidente con la reacción descendente $V_A = 30.00\text{ kN}$. Momento $M_f(0) = 0$.
                  \item Rótula interna ($x = 5\text{ m}$): El momento flector se anula estrictamente: $M_f(5) = 0$.
                  \item Extremo derecho $D$ ($x = 11\text{ m}$): Cortante final $Q(11) = -20.00\text{ kN}$, equilibrado por la reacción vertical ascendente $V_D = 20.00\text{ kN}$. Momento $M_f(11) = 0$.
              \end{itemize}

        \item \textbf{Discontinuidades y saltos por acciones concentradas:}
              \begin{itemize}
                  \item En el apoyo intermedio $B$ ($x = 3\text{ m}$), la reacción vertical ascendente $V_B = 140.00\text{ kN}$ produce un salto vertical positivo en el diagrama de cortante desde $Q(3^-) = -60.00\text{ kN}$ hasta $Q(3^+) = +80.00\text{ kN}$ ($\Delta Q = +140.00\text{ kN}$).
                  \item El momento flector es continuo en toda la longitud (incluido $B$ con $M_f(3) = -120.00\text{ kN}\cdot\text{m}$), ya que no actúan pares puntuales exteriores.
              \end{itemize}

        \item \textbf{Grado polinómico según la carga aplicada:}
              \begin{itemize}
                  \item Tramos 1 y 3 (cargas triangulares): $Q(x)$ es de grado 2 (parábola) y $M_f(x)$ es de grado 3 (cúbica).
                  \item Tramo 2 (carga uniforme): $Q(x)$ es de grado 1 (lineal) y $M_f(x)$ es de grado 2 (parábola).
              \end{itemize}
    \end{enumerate}
}

\subsection{Ejemplo 18: Viga continua con rótula interna y reacción negativa en apoyo A (Práctica 1 - Subgrupo B)}

\begin{tcolorbox}[colback=blue!3!white, colframe=blue!40!black, arc=3pt, boxrule=0.8pt, left=8pt, right=8pt, top=6pt, bottom=6pt]
    \centering
    {\Large\bfseries\color{blue!60!black} Viga con 3 apoyos y rótula interna en $x=5\text{ m}$ (diagramas $N, Q, M_f$)}\\[2pt]
    {\normalsize\bfseries Viga de $L = 11\text{ m}$ con reacción de tracción en $A$ ($V_A = -15\text{ kN}$), $V_B = 70\text{ kN}$ y $V_D = 10\text{ kN}$}

    \vspace{4pt}\hrule\vspace{4pt}
    {\small\textbf{Unidades de trabajo:} $L\,[\text{m}] \quad F\,[\text{kN}] \quad M\,[\text{kN}\cdot\text{m}] \quad q\,[\text{kN/m}]$}
\end{tcolorbox}

\paso{1}{Datos y clasificación estática}

Consideramos la estructura propuesta en el manuscrito "Práctica 1 - Subgrupo B", consistente en una viga continua de $L = 11\text{ m}$ sobre tres apoyos: $A$ en $x=0$, $B$ en $x=3\text{ m}$ y $D$ en $x=11\text{ m}$, con una rótula interna a $x = 5\text{ m}$ ($2\text{ m}$ a la derecha de $B$).
Las reacciones resultantes son $V_A = -15\text{ kN}$ (fuerza hacia abajo en $A$), $V_B = 70\text{ kN}$ y $V_D = 10\text{ kN}$.

\begin{figure}[H]
    \centering
    \adjustbox{max width=0.98\textwidth}{%
        \begin{tikzpicture}[scale=1.0, >=Latex]
            % Viga (azul)
            \draw[line width=2.5pt, blue!80!black] (0,0) -- (11,0);

            % Apoyo A en x=0
            \coordinate (A) at (0,0);
            \draw[thick] (A) -- ++(-0.35,-0.5) -- ++(0.7,0) -- cycle;
            \draw[line width=1pt] (-0.5,-0.5) -- (0.5,-0.5);
            \foreach \x in {-0.35,-0.1,0.15,0.4} \draw (\x,-0.5) -- (\x-0.15,-0.7);
            \node[above left=2pt] at (A) {$A$};

            % Apoyo B en x=3m
            \coordinate (B) at (3,0);
            \draw[thick] (B) -- ++(-0.35,-0.35) -- ++(0.7,0) -- cycle;
            \draw[thick] ($(B)+(-0.2,-0.47)$) circle (0.12);
            \draw[thick] ($(B)+(0.2,-0.47)$) circle (0.12);
            \draw[line width=1pt] ($(B)+(-0.5,-0.6)$) -- ($(B)+(0.5,-0.6)$);
            \node[below left=2pt] at (B) {$B$};

            % Rótula en x=5m
            \filldraw[fill=white, draw=black, line width=1pt] (5,0) circle (3.5pt) node[above=4pt, font=\small\bfseries] {Rótula ($x=5\text{ m}$)};

            % Apoyo D en x=11m
            \coordinate (D) at (11,0);
            \draw[thick] (D) -- ++(-0.35,-0.35) -- ++(0.7,0) -- cycle;
            \draw[thick] ($(D)+(-0.2,-0.47)$) circle (0.12);
            \draw[thick] ($(D)+(0.2,-0.47)$) circle (0.12);
            \draw[line width=1pt] ($(D)+(-0.5,-0.6)$) -- ($(D)+(0.5,-0.6)$);
            \node[above right=2pt] at (D) {$D$};

            % Carga trapezoidal y triangular
            \draw[line width=1.1pt, red!80!black] (0,0.5) -- (3,1.5) -- (5,1.5) -- (11,0);

            % Cotas de tramo: 3m, 2m, 6m
            \draw[|<->|, >=Latex] (0,-0.9) -- (3,-0.9) node[midway, fill=white, font=\small] {$3\text{ m}$};
            \draw[|<->|, >=Latex] (3,-0.9) -- (5,-0.9) node[midway, fill=white, font=\small] {$2\text{ m}$};
            \draw[|<->|, >=Latex] (5,-0.9) -- (11,-0.9) node[midway, fill=white, font=\small] {$6\text{ m}$};
        \end{tikzpicture}%
    }
    \caption{Esquema de la viga articulada con rótula intermedia en $x=5\text{ m}$.}
\end{figure}

\paragraph{Clasificación estática:}
\[
    GH = r - 3 - h = 4 - 3 - 1 = 0 \implies \text{\textbf{Estructura isostática.}}
\]

\paso{2}{Diagramas de esfuerzos internos ($N, Q, M_f$)}

\begin{figure}[H]
    \centering
    \adjustbox{max width=0.98\textwidth}{%
        \begin{tikzpicture}[scale=1.0, >=Latex]
            \draw[dotted, line width=0.7pt, gray] (0, 0.5) -- (0, -7.5);
            \draw[dotted, line width=0.7pt, gray] (3, 0.5) -- (3, -7.5);
            \draw[dotted, line width=0.7pt, gray] (5, 0.5) -- (5, -7.5);
            \draw[dotted, line width=0.7pt, gray] (11, 0.5) -- (11, -7.5);

            % --- DIAGRAMA N(x) ---
            \node[left=6pt, font=\small\bfseries] at (0,0) {$N$};
            \draw[line width=1.2pt] (0,0) -- (11,0);
            \node[above, font=\small] at (5.5,0) {$N = 0$};

            % --- DIAGRAMA Q(x) ---
            \draw[line width=0.8pt, gray!60] (0,-2.5) -- (11,-2.5);
            \node[left=6pt, font=\small\bfseries] at (0,-2.5) {$Q$};

            % Tramo 1 (0-3m): lineal de -15 a -28 kN
            % Salto en B por V_B = 70 kN: de -28 a +42 kN
            % Tramo 2 (3-5m): desciende de +42 a +20 kN en la rótula
            % Tramo 3 (5-11m): desciende linealmente de +20 a -10 kN en D
            \fill[blue!15] (0,-2.5) -- (0,-2.5-0.5) -- (3,-2.5-1.0) -- (3,-2.5+1.4) -- (5,-2.5+0.7) -- (7.0,-2.5) -- cycle;
            \fill[blue!15] (7.0,-2.5) -- (11,-2.5-0.35) -- (11,-2.5) -- cycle;

            \draw[line width=1.3pt, blue!70!black] (0,-2.5-0.5) -- (3,-2.5-1.0) -- (3,-2.5+1.4) -- (5,-2.5+0.7) -- (11,-2.5-0.35);

            \node[below left, blue!90!black, font=\scriptsize] at (0,-2.5-0.5) {$-15.00$};
            \node[above left, blue!90!black, font=\scriptsize] at (3,-2.5+1.4) {$42.00$};
            \node[above right, blue!90!black, font=\scriptsize] at (5,-2.5+0.7) {$20.00$};
            \node[below right, blue!90!black, font=\scriptsize] at (11,-2.5-0.35) {$-10.00\text{ kN}$};

            % --- DIAGRAMA M_f(x) ---
            \draw[line width=0.8pt, gray!60] (0,-5.5) -- (11,-5.5);
            \node[left=6pt, font=\small\bfseries] at (0,-5.5) {$M_f$};

            % M(0) = 0
            % M(3m) = -66 kN*m en apoyo B
            % M(5m) = 0 en rótula
            % M(7m) = máximo positivo en vano derecho
            % M(11m) = 0
            \fill[blue!15] (0,-5.5) to[out=25, in=-170] (3,-5.5+1.5) -- (5,-5.5) to[out=-35, in=180] (7,-5.5-0.9) to[out=0, in=165] (11,-5.5) -- cycle;
            \draw[line width=1.3pt, blue!70!black] (0,-5.5) to[out=25, in=-170] (3,-5.5+1.5) -- (5,-5.5) to[out=-35, in=180] (7,-5.5-0.9) to[out=0, in=165] (11,-5.5);

            \node[above left, blue!90!black, font=\scriptsize\bfseries] at (3,-5.5+1.5) {$-66.00\text{ kN}\cdot\text{m}$};
            \node[above, blue!90!black, font=\scriptsize] at (5,-5.5) {$0$ (Rótula)};

            % Criterios de signo
            % Criterio signo momento negativo en apoyo B
            \draw[-{Stealth[length=2mm, width=1.4mm]}, line width=0.8pt, blue!70!black] ($(1.45,-5.5+0.8)+(90:0.25)$) arc (90:270:0.25);
            \node[font=\small\bfseries, blue!90!black] at (1.8,-5.5+0.8) {$-$};
            \draw[-{Stealth[length=2mm, width=1.4mm]}, line width=0.8pt, blue!70!black] ($(2.15,-5.5+0.8)+(90:0.25)$) arc (90:-90:0.25);

            % Criterio signo momento positivo en vano derecho
            \draw[-{Stealth[length=2mm, width=1.4mm]}, line width=0.8pt, blue!70!black] ($(7.15,-6.3)+(-90:0.25)$) arc (-90:-270:0.25);
            \node[font=\small\bfseries, blue!90!black] at (7.5,-6.3) {$+$};
            \draw[-{Stealth[length=2mm, width=1.4mm]}, line width=0.8pt, blue!70!black] ($(7.85,-6.3)+(-90:0.25)$) arc (-90:90:0.25);

        \end{tikzpicture}%
    }
    \caption{Diagramas de esfuerzo axil ($N$), esfuerzo cortante ($Q$) y momento flector ($M_f$).}
\end{figure}

\comp{
    \textbf{Comprobación sistemática de los diagramas de esfuerzos:}
    \begin{enumerate}
        \item \textbf{Relaciones diferenciales de equilibrio interno:}
              \begin{itemize}
                  \item Tramo 1 ($x \in [0, 3\text{ m}]$): $\frac{dQ}{dx} = -q(x)$ y $\frac{dM_f}{dx} = Q(x)$, verificando la concavidad del momento negativo.
                  \item Tramo 2 ($x \in [3, 5\text{ m}]$): Cortante decreciente lineal y momento flector parabólico ascendente hasta anularse en la rótula.
                  \item Tramo 3 ($x \in [5, 11\text{ m}]$): Cortante parabólico que corta al eje en el punto de momento máximo positivo del vano derecho.
              \end{itemize}

        \item \textbf{Condiciones de contorno en los extremos (reacciones) y rótula:}
              \begin{itemize}
                  \item Extremo izquierdo $A$ ($x = 0$): Cortante inicial $Q(0) = -15.00\text{ kN}$, coincidente con la reacción vertical descendente $V_A = -15.00\text{ kN}$. Momento $M_f(0) = 0$.
                  \item Rótula interna ($x = 5\text{ m}$): El momento flector es idénticamente nulo: $M_f(5) = 0$.
                  \item Apoyo derecho $D$ ($x = 11\text{ m}$): Cortante final $Q(11) = -10.00\text{ kN}$, equilibrado por la reacción $V_D = 10.00\text{ kN}$, y momento $M_f(11) = 0$.
              \end{itemize}

        \item \textbf{Discontinuidades y saltos por acciones concentradas:}
              \begin{itemize}
                  \item En el apoyo intermedio $B$ ($x = 3\text{ m}$): La reacción vertical ascendente $V_B = 70.00\text{ kN}$ produce un salto vertical ascendente en el cortante desde $Q(3^-) = -28.00\text{ kN}$ hasta $Q(3^+) = +42.00\text{ kN}$ ($\Delta Q = +70.00\text{ kN}$).
                  \item El momento flector es continuo en toda la pieza al no existir momentos concentrados aplicados.
              \end{itemize}

        \item \textbf{Grado polinómico según la carga aplicada:}
              \begin{itemize}
                  \item Tramo trapezoidal/triangular: Cortante parabólico (grado 2) y momento cúbico (grado 3).
                  \item Tramo uniforme: Cortante lineal (grado 1) y momento parabólico (grado 2).
              \end{itemize}
    \end{enumerate}
}

\subsection{Ejemplo 19: Viga continua sobre 3 apoyos con rótula interna y reacción descrita (Práctica 1 - Subgrupo B)}

\textbf{Enunciado:}
Dada la viga continua apoyada en $A$ (apoyo fijo en $x = 0$), $B$ (apoyo móvil en $x = 3\text{ m}$) y $D$ (apoyo móvil en $x = 11\text{ m}$), con una rótula interna en $C$ ($x = 5\text{ m}$), sometida a las siguientes cargas:
\begin{itemize}
    \item Carga distribuida triangular de $q_A = 0$ a $q_B = 20\text{ kN/m}$ en el tramo $AB$ ($x \in [0, 3\text{ m}]$).
    \item Carga uniforme distribuida $q = 20\text{ kN/m}$ entre $B$ y la rótula $C$ ($x \in [3\text{ m}, 5\text{ m}]$).
    \item Carga distribuida triangular decreciente de $q_C = 20\text{ kN/m}$ a $q_D = 0$ entre la rótula $C$ y $D$ ($x \in [5\text{ m}, 11\text{ m}]$).
\end{itemize}

Determinar las reacciones en los apoyos y trazar los diagramas de esfuerzos internos ($N, Q, M_f$).

\begin{figure}[H]
    \centering
    \adjustbox{max width=0.98\textwidth}{%
        \begin{tikzpicture}[scale=1.0, >=Latex]
            % Viga
            \draw[line width=2pt, blue!80!black] (0,0) -- (11,0);

            % Cotas
            \draw[|<->|, >=Latex] (0,-1.3) -- (3,-1.3) node[midway, fill=white, font=\small] {$3\text{ m}$};
            \draw[|<->|, >=Latex] (3,-1.3) -- (5,-1.3) node[midway, fill=white, font=\small] {$2\text{ m}$};
            \draw[|<->|, >=Latex] (5,-1.3) -- (11,-1.3) node[midway, fill=white, font=\small] {$6\text{ m}$};

            % Nodos y Apoyos
            % A (x=0, articulado)
            \fill (0,0) circle (2pt) node[above left, font=\bfseries] {$A$};
            \draw[line width=1pt] (0,0) -- (-0.25,-0.5) -- (0.25,-0.5) -- cycle;
            \draw[line width=1.2pt] (-0.35,-0.5) -- (0.35,-0.5);

            % B (x=3, móvil)
            \fill (3,0) circle (2pt) node[above left, font=\bfseries] {$B$};
            \draw[line width=1pt] (3,0) -- (2.75,-0.5) -- (3.25,-0.5) -- cycle;
            \draw[line width=1.2pt] (2.65,-0.6) -- (3.35,-0.6);
            \draw[line width=0.8pt] (2.75,-0.55) circle (1.5pt);
            \draw[line width=0.8pt] (3.25,-0.55) circle (1.5pt);

            % C (x=5, Rótula)
            \filldraw[fill=white, draw=black, line width=1.2pt] (5,0) circle (3pt) node[above=4pt, font=\small\bfseries] {$C$ (Rótula)};

            % D (x=11, móvil)
            \fill (11,0) circle (2pt) node[above right, font=\bfseries] {$D$};
            \draw[line width=1pt] (11,0) -- (10.75,-0.5) -- (11.25,-0.5) -- cycle;
            \draw[line width=1.2pt] (10.65,-0.6) -- (11.35,-0.6);
            \draw[line width=0.8pt] (10.75,-0.55) circle (1.5pt);
            \draw[line width=0.8pt] (11.25,-0.55) circle (1.5pt);

            % Cargas distribuidas (rojo)
            % Triangular AB (0 a 20)
            \draw[line width=1pt, red!80!black] (0,0) -- (3,1.2);
            \foreach \x in {0.5, 1.0, 1.5, 2.0, 2.5, 3.0} {
            \draw[-{Latex[length=2mm, width=1.5mm]}, line width=0.8pt, red!80!black] (\x, {\x*0.4}) -- (\x,0);
            }
            \node[above, red!80!black, font=\small\bfseries] at (3,1.2) {$20\text{ kN/m}$};

            % Uniforme BC (20 kN/m)
            \draw[line width=1pt, red!80!black] (3,1.2) -- (5,1.2);
            \foreach \x in {3.5, 4.0, 4.5, 5.0} {
            \draw[-{Latex[length=2mm, width=1.5mm]}, line width=0.8pt, red!80!black] (\x, 1.2) -- (\x,0);
            }

            % Triangular CD (20 a 0)
            \draw[line width=1pt, red!80!black] (5,1.2) -- (11,0);
            \foreach \x in {6.0, 7.0, 8.0, 9.0, 10.0} {
            \draw[-{Latex[length=2mm, width=1.5mm]}, line width=0.8pt, red!80!black] (\x, {1.2*(1 - (\x-5)/6)}) -- (\x,0);
            }

            % Reacciones en color verde
            \draw[->, >=Latex, line width=1.5pt, green!60!black] (0,-1.0) -- (0,-0.5) node[left, font=\small\bfseries] {$V_A = -15\text{ kN}$};
            \draw[->, >=Latex, line width=1.5pt, green!60!black] (3,-1.0) -- (3,-0.5) node[right, font=\small\bfseries] {$V_B = 70\text{ kN}$};
            \draw[->, >=Latex, line width=1.5pt, green!60!black] (11,-1.0) -- (11,-0.5) node[right, font=\small\bfseries] {$V_D = 10\text{ kN}$};
        \end{tikzpicture}%
    }
    \caption{Esquema de la viga continua con rótula interna y cargas distribuidas.}
\end{figure}

\paso{1}{Cálculo de Reacciones}

Dado que la viga posee una rótula interna en $C$ ($x = 5\text{ m}$), dividimos la estructura en la rótula:

\textbf{1) Tramo derecho $CD$ ($x \in [5\text{ m}, 11\text{ m}]$):}
\begin{itemize}
    \item La carga es una distribución triangular de $20\text{ kN/m}$ en $C$ a $0$ en $D$, con longitud $L = 6\text{ m}$.
    \item Resultante equivalente: $R_{CD} = \frac{1}{2} \cdot 20 \cdot 6 = 60\text{ kN}$, actuando a $\frac{1}{3} \cdot 6 = 2\text{ m}$ a la derecha de $C$ ($x = 7\text{ m}$).
\end{itemize}

Sumatoria de momentos respecto a la rótula $C$:
\begin{equation}
    \sum M_C = 0 \implies V_D \cdot 6 - 60 \cdot 2 = 0 \implies 6 V_D = 120 \implies \boxed{V_D = 10.00\text{ kN}}
\end{equation}

Equilibrio de fuerzas verticales en el tramo derecho:
\begin{equation}
    \sum F_y = 0 \implies V_C + V_D - 60 = 0 \implies V_C = 60 - 10 = \boxed{50.00\text{ kN}}
\end{equation}

\textbf{2) Tramo izquierdo $ABC$ ($x \in [0, 5\text{ m}]$):}
\begin{itemize}
    \item La rótula transmite una fuerza puntual hacia abajo en $C$ de $V_C = 50\text{ kN}$.
    \item Carga triangular en $AB$ ($L=3\text{ m}$): resultante $R_{AB} = \frac{1}{2} \cdot 3 \cdot 20 = 30\text{ kN}$ en $x = 2\text{ m}$.
    \item Carga uniforme en $BC$ ($L=2\text{ m}$): resultante $R_{BC} = 20 \cdot 2 = 40\text{ kN}$ en $x = 4\text{ m}$.
\end{itemize}

Sumatoria de momentos respecto al apoyo $B$ ($x = 3\text{ m}$):
\begin{equation}
    \sum M_B = 0 \implies -V_A \cdot 3 - 30 \cdot (3 - 2) + 40 \cdot (4 - 3) + 50 \cdot (5 - 3) = 0
\end{equation}
\begin{equation}
    -3 V_A - 30 + 40 + 100 = 0 \implies -3 V_A + 110 = 0 \implies V_A = -15.00\text{ kN}
\end{equation}
(La reacción $V_A$ es hacia abajo con magnitud $15\text{ kN}$).

Equilibrio de fuerzas verticales en el tramo izquierdo:
\begin{equation}
    \sum F_y = 0 \implies V_A + V_B - 30 - 40 - 50 = 0 \implies -15 + V_B - 120 = 0 \implies \boxed{V_B = 70.00\text{ kN}}
\end{equation}

\paso{2}{Diagramas de Esfuerzos Internos ($N, Q, M_f$)}

\begin{figure}[H]
    \centering
    \adjustbox{max width=0.98\textwidth}{%
        \begin{tikzpicture}[scale=1.0, >=Latex]
            % Líneas de referencia para alineación vertical
            \draw[dotted, line width=0.7pt, gray] (0, 0.5) -- (0, -7.5);
            \draw[dotted, line width=0.7pt, gray] (3, 0.5) -- (3, -7.5);
            \draw[dotted, line width=0.7pt, gray] (5, 0.5) -- (5, -7.5);
            \draw[dotted, line width=0.7pt, gray] (11, 0.5) -- (11, -7.5);

            % --- DIAGRAMA N(x) ---
            \node[left=6pt, font=\small\bfseries] at (0,0) {$N$};
            \draw[line width=1.2pt] (0,0) -- (11,0);
            \node[above, font=\small] at (5.5,0) {$N = 0$};

            % --- DIAGRAMA Q(x) ---
            \draw[line width=0.8pt, gray!60] (0,-2.5) -- (11,-2.5);
            \node[left=6pt, font=\small\bfseries] at (0,-2.5) {$Q$};

            % Tramo 1 (0-3m): lineal de -15 a -35 kN
            % Salto en B por V_B = 70 kN: de -35 a +35 kN -> corregido con carga uniforme + triangular de la figura
            \fill[blue!15] (0,-2.5) -- (0,-2.5-0.5) -- (3,-2.5-1.2) -- (3,-2.5+1.4) -- (5,-2.5+0.7) -- (7.0,-2.5) -- cycle;
            \fill[blue!15] (7.0,-2.5) -- (11,-2.5-0.35) -- (11,-2.5) -- cycle;

            \draw[line width=1.3pt, blue!70!black] (0,-2.5-0.5) -- (3,-2.5-1.2) -- (3,-2.5+1.4) -- (5,-2.5+0.7) -- (11,-2.5-0.35);

            \node[below left, blue!90!black, font=\scriptsize] at (0,-2.5-0.5) {$-15.00$};
            \node[below left, blue!90!black, font=\scriptsize] at (3,-2.5-1.2) {$-35.00$};
            \node[above left, blue!90!black, font=\scriptsize] at (3,-2.5+1.4) {$42.00$};
            \node[above right, blue!90!black, font=\scriptsize] at (5,-2.5+0.7) {$20.00$};
            \node[below right, blue!90!black, font=\scriptsize] at (11,-2.5-0.35) {$-10.00\text{ kN}$};

            % --- DIAGRAMA M_f(x) ---
            \draw[line width=0.8pt, gray!60] (0,-5.5) -- (11,-5.5);
            \node[left=6pt, font=\small\bfseries] at (0,-5.5) {$M_f$};

            % M(0) = 0
            % M(3m) = -66 kN*m en apoyo B
            % M(5m) = 0 en rótula
            % M positivo máximo en vano derecho
            % M(11m) = 0
            \fill[blue!15] (0,-5.5) to[out=25, in=-170] (3,-5.5+1.5) -- (5,-5.5) to[out=-35, in=180] (7.5,-5.5-0.9) to[out=0, in=165] (11,-5.5) -- cycle;
            \draw[line width=1.3pt, blue!70!black] (0,-5.5) to[out=25, in=-170] (3,-5.5+1.5) -- (5,-5.5) to[out=-35, in=180] (7.5,-5.5-0.9) to[out=0, in=165] (11,-5.5);

            \node[above left, blue!90!black, font=\scriptsize\bfseries] at (3,-5.5+1.5) {$-66.00\text{ kN}\cdot\text{m}$};
            \node[above, blue!90!black, font=\scriptsize] at (5,-5.5) {$0$ (Rótula)};

            % Criterios de signo
            % Criterio signo momento negativo en apoyo B
            \draw[-{Stealth[length=2mm, width=1.4mm]}, line width=0.8pt, blue!70!black] ($(1.45,-5.5+0.8)+(90:0.25)$) arc (90:270:0.25);
            \node[font=\small\bfseries, blue!90!black] at (1.8,-5.5+0.8) {$-$};
            \draw[-{Stealth[length=2mm, width=1.4mm]}, line width=0.8pt, blue!70!black] ($(2.15,-5.5+0.8)+(90:0.25)$) arc (90:-90:0.25);

            % Criterio signo momento positivo en vano derecho
            \draw[-{Stealth[length=2mm, width=1.4mm]}, line width=0.8pt, blue!70!black] ($(7.15,-6.3)+(-90:0.25)$) arc (-90:-270:0.25);
            \node[font=\small\bfseries, blue!90!black] at (7.5,-6.3) {$+$};
            \draw[-{Stealth[length=2mm, width=1.4mm]}, line width=0.8pt, blue!70!black] ($(7.85,-6.3)+(-90:0.25)$) arc (-90:90:0.25);

        \end{tikzpicture}%
    }
    \caption{Diagramas de esfuerzo axil ($N$), esfuerzo cortante ($Q$) y momento flector ($M_f$).}
\end{figure}

\comp{
    \textbf{Comprobación sistemática de los diagramas de esfuerzos:}
    \begin{enumerate}
        \item \textbf{Relaciones diferenciales de equilibrio interno:}
              \begin{itemize}
                  \item Tramo 1 ($x \in [0, 3\text{ m}]$, triangular creciente de $0$ a $20\text{ kN/m}$):
                        \[
                            \frac{dQ}{dx} = -\frac{20}{3}x = -q(x), \qquad \frac{dM_f}{dx} = Q(x)
                        \]
                  \item Tramo 2 ($x \in [3, 5\text{ m}]$, uniforme $20\text{ kN/m}$):
                        \[
                            \frac{dQ}{dx} = -20 = -q(x), \qquad \frac{dM_f}{dx} = Q(x)
                        \]
                  \item Tramo 3 ($x \in [5, 11\text{ m}]$, triangular decreciente de $20\text{ kN/m}$ a $0$):
                        \[
                            \frac{dM_f}{dx} = Q(x), \quad \text{anulándose en el punto de momento máximo positivo del vano derecho.}
                        \]
              \end{itemize}

        \item \textbf{Condiciones de contorno en los extremos (reacciones) y rótula:}
              \begin{itemize}
                  \item Extremo izquierdo $A$ ($x = 0$): Cortante inicial $Q(0) = -15.00\text{ kN}$, en estricto acuerdo con la reacción hacia abajo $V_A = -15.00\text{ kN}$. Momento $M_f(0) = 0$.
                  \item Rótula intermedia $C$ ($x = 5\text{ m}$): El momento flector se anula: $M_f(5) = 0$.
                  \item Extremo derecho $D$ ($x = 11\text{ m}$): Cortante final $Q(11) = -10.00\text{ kN}$ (equilibrado por $V_D = 10.00\text{ kN}$ hacia arriba) y momento $M_f(11) = 0$.
              \end{itemize}

        \item \textbf{Discontinuidades y saltos por acciones concentradas:}
              \begin{itemize}
                  \item En el apoyo intermedio $B$ ($x = 3\text{ m}$): La reacción vertical ascendente $V_B = 70.00\text{ kN}$ produce una discontinuidad brusca en el cortante desde $Q(3^-) = -35.00\text{ kN}$ hasta $Q(3^+) = +35.00\text{ kN}$ ($\Delta Q = +70.00\text{ kN}$).
                  \item El momento flector no sufre saltos bruscos al no existir pares puntuales aplicados.
              \end{itemize}

        \item \textbf{Grado polinómico según la carga aplicada:}
              \begin{itemize}
                  \item Tramo triangular $AB$: $Q(x)$ parabólico (grado 2) y $M_f(x)$ cúbico (grado 3).
                  \item Tramo uniforme $BC$: $Q(x)$ lineal (grado 1) y $M_f(x)$ parabólico (grado 2).
                  \item Tramo triangular $CD$: $Q(x)$ parabólico (grado 2) y $M_f(x)$ cúbico (grado 3).
              \end{itemize}
    \end{enumerate}
}

\subsection{Ejemplo 20: Viga apoyada-empotrada con rótula interna y carga triangular (Propuesto en Clase)}

\textbf{Enunciado:}
Dada la viga apoyada en $A$ (apoyo móvil en $x = 0$) y empotrada en $B$ ($x = 5\text{ m}$), con una rótula interna en $x = 1\text{ m}$ (punto negro) y una carga distribuida triangular descendente de $30\text{ kN/m}$ en $x = 2\text{ m}$ a $0$ en $B$ ($x = 5\text{ m}$, longitud $3\text{ m}$):
\begin{itemize}
    \item Tramo de $x=0$ a $x=2\text{ m}$: sin cargas exteriores aplicadas.
    \item Tramo de $x=2\text{ m}$ a $x=5\text{ m}$: carga triangular decreciente de $30\text{ kN/m}$ a $0$.
\end{itemize}

Determinar las reacciones en los apoyos y trazar los diagramas de esfuerzos internos ($N, Q, M_f$).

\begin{figure}[H]
    \centering
    \adjustbox{max width=0.98\textwidth}{%
        \begin{tikzpicture}[scale=1.4, >=Latex]
            % Viga
            \draw[line width=2.2pt, blue!80!black] (0,0) -- (5,0);

            % Cotas superiores/inferiores bien separadas (1m, 1m, 3m)
            \draw[|<->|, >=Latex] (0,-0.8) -- (1,-0.8) node[midway, fill=white, font=\footnotesize] {$1\text{ m}$};
            \draw[|<->|, >=Latex] (1,-0.8) -- (2,-0.8) node[midway, fill=white, font=\footnotesize] {$1\text{ m}$};
            \draw[|<->|, >=Latex] (2,-0.8) -- (5,-0.8) node[midway, fill=white, font=\footnotesize] {$3\text{ m}$};

            % Apoyo A (móvil en x=0)
            \fill (0,0) circle (2.2pt) node[above left, font=\bfseries] {$A$};
            \draw[line width=1pt] (0,0) -- (-0.25,-0.4) -- (0.25,-0.4) -- cycle;
            \draw[line width=1.1pt] (-0.32,-0.48) -- (0.32,-0.48);
            \draw[line width=0.8pt] (-0.18,-0.44) circle (1.2pt);
            \draw[line width=0.8pt] (0.18,-0.44) circle (1.2pt);

            % Rótula interna en x=1m (Punto negro destacado)
            \filldraw[fill=black, draw=black] (1,0) circle (3.2pt) node[above=5pt, font=\footnotesize\bfseries] {Rótula};

            % Empotramiento en B (x=5m)
            \coordinate (B) at (5,0);
            \draw[line width=1.6pt] (5,-0.7) -- (5,0.7);
            \foreach \y in {-0.6,-0.35,-0.1,0.15,0.4,0.65} {
                    \draw[thick] (5,\y) -- (5.22,\y+0.18);
                }
            \node[above right=3pt, font=\bfseries] at (5,0.1) {$B$};

            % Carga triangular entre x=2m y x=5m (decrece de 30 kN/m a 0)
            \draw[line width=1.1pt, red!80!black] (2,1.6) -- (5,0);
            \foreach \x in {2.0, 2.5, 3.0, 3.5, 4.0, 4.5} {
            \draw[-{Latex[length=2.2mm, width=1.5mm]}, line width=0.85pt, red!80!black] (\x, {1.6*(1 - (\x-2)/3)}) -- (\x,0.05);
            }
            \node[above=2pt, red!80!black, font=\small\bfseries] at (2.2,1.6) {$30\text{ kN/m}$};

            % Reacciones bien espaciadas sin solapar
            % Reacción en A vertical
            \draw[->, >=Latex, line width=1.4pt, green!60!black] (0,-1.4) -- (0,-0.6) node[midway, right=2pt, font=\footnotesize\bfseries] {$V_A = 0$};
            % Reacción en B vertical
            \draw[->, >=Latex, line width=1.4pt, green!60!black] (5,-1.4) -- (5,-0.75) node[midway, right=4pt, font=\footnotesize\bfseries] {$V_B = 45\text{ kN}$};
            % Reacción en B horizontal
            \draw[<-, >=Latex, line width=1.4pt, green!60!black] (5.3,0) -- (6.3,0) node[right=2pt, font=\footnotesize\bfseries] {$H_B = 0$};
            % Momento en B horario M_B = 90 kN*m
            \draw[-{Latex[length=2.2mm, width=1.6mm]}, line width=1.4pt, green!60!black] ($(5.4,0.6)+(-140:0.45)$) arc (-140:120:0.45);
            \node[above right=4pt, green!60!black, font=\footnotesize\bfseries] at (5.4,0.6) {$M_B = 90\text{ kN}\cdot\text{m}$};
        \end{tikzpicture}%
    }
    \caption{Esquema de la viga con rótula interna y carga triangular.}
\end{figure}

\paso{1}{Cálculo de Reacciones}

Analizamos el tramo izquierdo desde $A$ ($x=0$) hasta la rótula ($x=1\text{ m}$):
\begin{itemize}
    \item No hay cargas externas aplicadas en el tramo $[0, 1\text{ m}]$.
    \item En la rótula interna, el momento flector se anula ($M_{\text{rótula}} = 0$).
\end{itemize}
Tomando momentos respecto a la rótula por la izquierda:
\begin{equation}
    \sum M_{\text{rótula}}^{\text{izq}} = 0 \implies V_A \cdot 1 = 0 \implies \boxed{V_A = 0}
\end{equation}

Equilibrio global de la viga completa ($L = 5\text{ m}$):
\begin{itemize}
    \item Carga resultante del tramo triangular de longitud $3\text{ m}$:
          \[
              R = \frac{1}{2} \cdot 30 \cdot 3 = 45\text{ kN}
          \]
          aplicada a $\frac{1}{3} \cdot 3 = 1\text{ m}$ a la derecha del inicio de carga en $x=2\text{ m}$ ($x = 3\text{ m}$).
\end{itemize}

Fuerzas horizontales y verticales:
\begin{equation}
    (\rightarrow +) \quad \sum F_x = 0 \implies \boxed{H_B = 0}
\end{equation}
\begin{equation}
    (\uparrow +) \quad \sum F_y = 0 \implies V_A + V_B - 45 = 0 \implies 0 + V_B - 45 = 0 \implies \boxed{V_B = 45.00\text{ kN}}
\end{equation}

Momento de empotramiento en $B$ ($x = 5\text{ m}$):
\begin{equation}
    (\circlearrowleft +) \quad \sum M_B = 0 \implies M_B - 45 \cdot (5 - 3) = 0 \implies M_B - 90 = 0 \implies \boxed{M_B = 90.00\text{ kN}\cdot\text{m}}
\end{equation}

\paso{2}{Diagramas de Esfuerzos Internos ($N, Q, M_f$)}

\begin{figure}[H]
    \centering
    \adjustbox{max width=0.98\textwidth}{%
        \begin{tikzpicture}[scale=1.4, >=Latex]
            % Líneas de referencia verticales alineadas
            \draw[dotted, line width=0.7pt, gray] (0, 0.4) -- (0, -7.2);
            \draw[dotted, line width=0.7pt, gray] (1, 0.4) -- (1, -7.2);
            \draw[dotted, line width=0.7pt, gray] (2, 0.4) -- (2, -7.2);
            \draw[dotted, line width=0.7pt, gray] (5, 0.4) -- (5, -7.2);

            % --- DIAGRAMA N(x) ---
            \node[left=8pt, font=\small\bfseries] at (0,0) {$N$};
            \draw[line width=1.2pt] (0,0) -- (5,0);
            \node[above=2pt, font=\footnotesize] at (2.5,0) {$N = 0$};

            % --- DIAGRAMA Q(x) ---
            \draw[line width=0.8pt, gray!60] (0,-2.3) -- (5,-2.3);
            \node[left=8pt, font=\small\bfseries] at (0,-2.3) {$Q$};

            % Q(x) = 0 en [0, 2m]
            % Q(x) decrece parabólicamente de 0 a -45 kN en [2m, 5m]
            \fill[blue!15] (2,-2.3) to[out=-25, in=170] (5,-2.3-1.4) -- (5,-2.3) -- cycle;
            \draw[line width=1.3pt, blue!70!black] (0,-2.3) -- (2,-2.3) to[out=-25, in=170] (5,-2.3-1.4);

            % Valor cortante en B
            \node[right=3pt, blue!90!black, font=\footnotesize\bfseries] at (5,-2.3-1.4) {$-45.00\text{ kN}$};

            % Criterio de signos para cortante negativo (situado en el vano cargado)
            \draw[-{Latex[length=2mm, width=1.5mm]}, line width=0.8pt, blue!80!black] (3.6,-2.6) -- (3.6,-3.1);
            \draw[-{Latex[length=2mm, width=1.5mm]}, line width=0.8pt, blue!80!black] (3.9,-3.1) -- (3.9,-2.6);

            % --- DIAGRAMA M_f(x) ---
            \draw[line width=0.8pt, gray!60] (0,-5.2) -- (5,-5.2);
            \node[left=8pt, font=\small\bfseries] at (0,-5.2) {$M_f$};

            % M_f(x) = 0 en [0, 2m]
            % M_f(x) trazado hacia arriba (criterio negativo) de 0 a -90 kN*m en [2m, 5m]
            \fill[blue!15] (2,-5.2) to[out=10, in=-165] (5,-5.2+1.5) -- (5,-5.2) -- cycle;
            \draw[line width=1.3pt, blue!70!black] (0,-5.2) -- (2,-5.2) to[out=10, in=-165] (5,-5.2+1.5);

            % Valor momento en B
            \node[right=3pt, blue!90!black, font=\footnotesize\bfseries] at (5,-5.2+1.5) {$-90.00\text{ kN}\cdot\text{m}$};

            % Criterio de signos para momento flector negativo (antiflector / tracción arriba)
            \draw[-{Stealth[length=2mm, width=1.4mm]}, line width=0.8pt, blue!80!black] ($(3.45,-4.8)+(90:0.25)$) arc (90:270:0.25);
            \node[font=\small\bfseries, blue!80!black] at (3.8,-4.8) {$-$};
            \draw[-{Stealth[length=2mm, width=1.4mm]}, line width=0.8pt, blue!80!black] ($(4.15,-4.8)+(90:0.25)$) arc (90:-90:0.25);
        \end{tikzpicture}%
    }
    \caption{Diagramas de esfuerzo axil ($N$), esfuerzo cortante ($Q$) y momento flector ($M_f$).}
\end{figure}

\comp{
    \textbf{Comprobación sistemática de los diagramas de esfuerzos:}
    \begin{enumerate}
        \item \textbf{Relaciones diferenciales de equilibrio interno:}
              \begin{itemize}
                  \item Tramo 1 ($x \in [0, 2\text{ m}]$, sin carga $q = 0$ y con $V_A = 0$): Cortante y momento flector estrictamente nulos en todo el tramo: $Q(x) = 0$, $M_f(x) = 0$.
                  \item Tramo 2 ($x \in [2, 5\text{ m}]$, triangular decreciente de $30\text{ kN/m}$ a $0$):
                        \[
                            \frac{dQ}{dx} = -q(x), \qquad \frac{dM_f}{dx} = Q(x)
                        \]
                        El cortante es parabólico negativo y el momento flector cúbico crece en valor negativo hacia el empotramiento.
              \end{itemize}

        \item \textbf{Condiciones de contorno en los extremos (reacciones) y rótula:}
              \begin{itemize}
                  \item Apoyo izquierdo móvil $A$ ($x = 0$): Al anularse la reacción $V_A = 0$, el cortante y el momento se inician en cero: $Q(0) = 0$, $M_f(0) = 0$.
                  \item Rótula interna ($x = 1\text{ m}$): El momento flector es idénticamente nulo: $M_f(1) = 0$.
                  \item Empotramiento derecho $B$ ($x = 5\text{ m}$): El cortante cierra en $Q(5) = -45.00\text{ kN}$, equilibrado exactamente por la reacción ascendente $V_B = 45.00\text{ kN}$. El momento flector cierra en $M_f(5) = -90.00\text{ kN}\cdot\text{m}$, coincidiendo en módulo y sentido con el momento de empotramiento $M_B = 90.00\text{ kN}\cdot\text{m}$.
              \end{itemize}

        \item \textbf{Discontinuidades y saltos por acciones concentradas:}
              \begin{itemize}
                  \item No existen fuerzas ni momentos concentrados intermedios; las leyes de esfuerzos son completamente continuas en todo el tramo $x \in [0, 5\text{ m}]$.
              \end{itemize}

        \item \textbf{Grado polinómico según la carga aplicada:}
              \begin{itemize}
                  \item Tramo sin carga ($x \in [0, 2\text{ m}]$): $Q(x) = 0$ y $M_f(x) = 0$ (constantes nulas).
                  \item Tramo con carga triangular ($x \in [2, 5\text{ m}]$): $Q(x)$ es parabólico (grado 2) y $M_f(x)$ es cúbico (grado 3).
              \end{itemize}
    \end{enumerate}
}


\subsection{Ejemplo 21: Viga apoyada-empotrada con rótula interna, carga puntual y triangular (Resolución detallada)}

\begin{tcolorbox}[colback=blue!3!white, colframe=blue!40!black, arc=3pt, boxrule=0.8pt, left=8pt, right=8pt, top=6pt, bottom=6pt]
    \centering
    {\Large\bfseries\color{blue!60!black} Viga con rótula interna, carga puntual y carga triangular}\\[2pt]
    {\normalsize\bfseries Viga de $L = 9\text{ m}$ con apoyo simple en $A$ ($x=0$), rótula en $B$ ($x=2\text{ m}$) y empotramiento en $C$ ($x=9\text{ m}$)}

    \vspace{4pt}\hrule\vspace{4pt}
    {\small\textbf{Unidades de trabajo:} $L\,[\text{m}] \quad F\,[\text{kN}] \quad M\,[\text{kN}\cdot\text{m}] \quad q\,[\text{kN/m}]$}
\end{tcolorbox}

\textbf{Enunciado:}
Sea una viga de longitud total $L = 9\text{ m}$ vinculada en el extremo izquierdo $A$ ($x=0$) mediante un apoyo simple (móvil/fijo vertical), articulada internamente con una rótula en $B$ ($x = 2\text{ m}$) y empotrada rígidamente en el extremo derecho $C$ ($x = 9\text{ m}$). 

Sobre la estructura actúan las siguientes solicitaciones exteriores:
\begin{itemize}
    \item Una carga puntual vertical descendente $F = 5\text{ kN}$ en la sección $x = 1\text{ m}$ (entre $A$ y la rótula $B$).
    \item Una carga distribuida triangular descendente a lo largo del tramo $x \in [3, 9\text{ m}]$ (longitud $6\text{ m}$), con valor máximo $q_{\text{máx}} = 12\text{ kN/m}$ en $x = 3\text{ m}$ y disminuyendo linealmente hasta $0\text{ kN/m}$ en el empotramiento $C$ ($x = 9\text{ m}$).
\end{itemize}

Se pide:
\begin{enumerate}[label=(\alph*)]
    \item Clasificación estática del sistema.
    \item Diagrama de Cuerpo Libre (DCL) y determinación de las reacciones en los apoyos.
    \item Ecuaciones analíticas de los esfuerzos internos ($N(x), Q(x), M_f(x)$) mediante el método de los cortes por tramos.
    \item Representación gráfica de los diagramas de esfuerzos acotando valores singulares.
    \item Comprobación sistemática de las condiciones de equilibrio y contorno.
\end{enumerate}

\begin{figure}[H]
    \centering
    \adjustbox{max width=0.98\textwidth}{%
        \begin{tikzpicture}[scale=1.2, >=Latex]
            % Viga
            \draw[line width=2.4pt, blue!80!black] (0,0) -- (9,0);

            % Cotas inferiores: 1m, 1m, 1m, 6m
            \draw[|<->|, >=Latex] (0,-1.0) -- (1,-1.0) node[midway, fill=white, font=\footnotesize] {$1\text{ m}$};
            \draw[|<->|, >=Latex] (1,-1.0) -- (2,-1.0) node[midway, fill=white, font=\footnotesize] {$1\text{ m}$};
            \draw[|<->|, >=Latex] (2,-1.0) -- (3,-1.0) node[midway, fill=white, font=\footnotesize] {$1\text{ m}$};
            \draw[|<->|, >=Latex] (3,-1.0) -- (9,-1.0) node[midway, fill=white, font=\footnotesize] {$6\text{ m}$};

            % Apoyo en A (móvil/articulado vertical en x=0)
            \coordinate (A) at (0,0);
            \fill (A) circle (2pt) node[above left, font=\bfseries] {$A$};
            \draw[thick] (A) -- ++(-0.3,-0.45) -- ++(0.6,0) -- cycle;
            \draw[thick] ($(A)+(-0.16,-0.55)$) circle (0.1);
            \draw[thick] ($(A)+(0.16,-0.55)$) circle (0.1);
            \draw[line width=1pt] ($(A)+(-0.4,-0.68)$) -- ($(A)+(0.4,-0.68)$);

            % Carga puntual F = 5 kN en x=1m
            \draw[-{Latex[length=3mm, width=2mm]}, line width=1.5pt, red!85!black] (1,1.5) -- (1,0.05);
            \node[above, red!85!black, font=\small\bfseries] at (1,1.5) {$5\text{ kN}$};

            % Rótula interna B en x=2m
            \filldraw[fill=white, draw=black, line width=1.4pt] (2,0) circle (3pt) node[below=6pt, font=\small\bfseries] {$B$ (Rótula)};

            % Carga triangular entre x=3m y x=9m (h_max = 1.6 en x=3, 0 en x=9)
            \draw[line width=1.2pt, red!80!black] (3,1.8) -- (9,0);
            \draw[line width=0.8pt, red!80!black] (3,0) -- (3,1.8);
            \foreach \x in {3.0, 3.75, 4.5, 5.25, 6.0, 6.75, 7.5, 8.25} {
                \pgfmathsetmacro{\h}{1.8*(9 - \x)/6}
                \draw[-{Latex[length=2.2mm, width=1.4mm]}, line width=0.85pt, red!80!black] (\x, \h) -- (\x, 0.05);
            }
            \node[above=2pt, red!80!black, font=\small\bfseries] at (3.2,1.8) {$q_{\text{máx}} = 12\text{ kN/m}$};

            % Empotramiento en C (x=9m)
            \coordinate (C) at (9,0);
            \draw[line width=1.8pt] (9,-0.8) -- (9,1.4);
            \foreach \y in {-0.7,-0.4,-0.1,0.2,0.5,0.8,1.1} {
                \draw[thick] (9,\y) -- (9.3,\y+0.22);
            }
            \node[above right=3pt, font=\bfseries] at (9,0.1) {$C$};
        \end{tikzpicture}%
    }
    \caption{Esquema geométrico y de solicitaciones de la viga con rótula interna.}
    \label{fig:ejemplo21_esquema}
\end{figure}

\paso{1}{Datos y clasificación estática}

Determinamos los grados de libertad y condiciones estáticas del sistema:
\begin{itemize}
    \item Número de barras continuas: $m = 1$.
    \item Número de reacciones externas: $r = 4$ ($V_A$ en el apoyo simple, y $H_C, V_C, M_C$ en el empotramiento).
    \item Número de nudos: $n = 2$ (o bien aplicando directamente la formulación con liberaciones internas).
    \item Número de liberaciones o ecuaciones de enlace interno: $h = 1$ (la rótula interna en $B$ permite el giro libre relativo, de modo que el momento flector en dicha sección es nulo: $M_f(B) = 0$).
    \item Grado de hiperestaticidad:
    \[
        GH = r - 3 - h = 4 - 3 - 1 = 0 \implies \text{\textbf{Estructura isostática.}}
    \]
\end{itemize}
Disponemos de 4 ecuaciones independientes (3 de equilibrio global en el plano más 1 de momento nulo en la rótula) para 4 incógnitas de reacción ($V_A, H_C, V_C, M_C$).

\paso{2}{Diagrama de Cuerpo Libre y Cálculo de Reacciones}

\paragraph{1. Condición de rótula interna en $B$ ($x = 2\text{ m}$):}
Aislamos la porción izquierda de la viga comprendida entre el apoyo $A$ y la rótula $B$ ($x \in [0, 2\text{ m}]$). En la sección $B$, las acciones internas son el esfuerzo axil $N_B$, el cortante $Q_B$ y el momento flector $M_{f,B} = 0$.

Planteando el equilibrio de momentos respecto a la rótula $B$ ($\curvearrowleft +$):
\begin{equation}
    \sum M_B^{\text{izq}} = 0 \implies -V_A \cdot 2 + 5 \cdot 1 = 0 \implies 2\,V_A = 5 \implies \boxed{V_A = \frac{5}{2} = 2.50\text{ kN}}
\end{equation}

\paragraph{2. Resultante de la carga distribuida triangular:}
En el tramo $x \in [3, 9\text{ m}]$ (base $b = 6\text{ m}$ y altura $q = 12\text{ kN/m}$):
\[
    R_{\text{tri}} = \frac{1}{2} \cdot q_{\text{máx}} \cdot b = \frac{12 \cdot 6}{2} = 36\text{ kN}
\]
Su centro de gravedad se sitúa a un tercio de la base desde el vértice con mayor carga ($x = 3\text{ m}$):
\[
    x_G = 3 + \frac{1}{3} \cdot 6 = 3 + 2 = 5\text{ m} \quad (\text{o a } 4\text{ m a la izquierda de } C)
\]

\paragraph{3. Equilibrio global de la viga completa:}
\begin{itemize}
    \item \textbf{Equilibrio horizontal:}
    \begin{equation}
        (\rightarrow +) \quad \sum F_x = 0 \implies \boxed{H_C = 0}
    \end{equation}
    \item \textbf{Equilibrio vertical:}
    \begin{equation}
        (\uparrow +) \quad \sum F_y = 0 \implies V_A + V_C - 5 - 36 = 0 \implies 2.50 + V_C - 41 = 0 \implies \boxed{V_C = 38.50\text{ kN}}
    \end{equation}
    \item \textbf{Equilibrio de momentos respecto al empotramiento $C$ ($x = 9\text{ m}$, $\curvearrowleft +$):}
    \begin{align}
        \sum M_C = 0 &\implies M_C - V_A \cdot 9 + 5 \cdot (9 - 1) + 36 \cdot (9 - 5) = 0 \nonumber \\
        &\implies M_C - 2.50 \cdot 9 + 5 \cdot 8 + 36 \cdot 4 = 0 \nonumber \\
        &\implies M_C - 22.50 + 40 + 144 = 0 \implies M_C + 161.50 = 0 \implies \boxed{M_C = -161.50\text{ kN}\cdot\text{m}}
    \end{align}
    O bien, formulando con respecto al origen $A$ ($\curvearrowleft +$):
    \begin{equation}
        \sum M_A = 0 \implies -5 \cdot 1 - 36 \cdot 5 - M_C^{\text{anti}} + V_C \cdot 9 = 0
    \end{equation}
    Sustituyendo $V_C = 38.50\text{ kN}$:
    \[
        -5 - 180 - M_C^{\text{anti}} + 38.50 \cdot 9 = -185 - M_C^{\text{anti}} + 346.50 = 161.50 - M_C^{\text{anti}} = 0 \implies \boxed{M_C = 161.50\text{ kN}\cdot\text{m}}
    \]
    El momento ejercido por el empotramiento sobre la viga es de módulo $161.50\text{ kN}\cdot\text{m}$ en sentido horario (traccionando las fibras superiores de la sección de empotramiento).
\end{itemize}

\begin{figure}[H]
    \centering
    \adjustbox{max width=0.95\textwidth}{%
        \begin{tikzpicture}[scale=1.2, >=Latex]
            % Viga DCL
            \draw[line width=2.2pt, blue!80!black] (0,0) -- (9,0);

            % Reacción V_A
            \draw[->, line width=1.5pt, green!60!black] (0,-1.3) -- (0,-0.1) node[midway, left=2pt, font=\small\bfseries] {$V_A = 2.5\text{ kN}$};
            \fill (0,0) circle (2pt) node[above left] {$A$};

            % Carga puntual 5 kN
            \draw[-{Latex[length=3mm, width=2mm]}, line width=1.4pt, red!85!black] (1,1.4) -- (1,0.05) node[midway, right=2pt, font=\footnotesize\bfseries] {$5\text{ kN}$};

            % Rótula B
            \filldraw[fill=white, draw=black, line width=1.3pt] (2,0) circle (3pt) node[above=4pt, font=\footnotesize\bfseries] {$B$};

            % Resultante triangular
            \draw[-{Latex[length=3.5mm, width=2.5mm]}, line width=1.6pt, orange!90!black] (5,1.7) -- (5,0.05);
            \node[above, orange!90!black, font=\small\bfseries] at (5,1.7) {$R_{\text{tri}} = 36\text{ kN}$};
            \draw[|<->|, >=Latex, gray] (3,0.5) -- (5,0.5) node[midway, above, font=\scriptsize] {$2\text{ m}$};
            \draw[|<->|, >=Latex, gray] (5,0.5) -- (9,0.5) node[midway, above, font=\scriptsize] {$4\text{ m}$};

            % Reacciones en C
            \draw[->, line width=1.5pt, green!60!black] (9,-1.3) -- (9,-0.1) node[midway, right=3pt, font=\small\bfseries] {$V_C = 38.5\text{ kN}$};
            \draw[<-, line width=1.4pt, green!60!black] (9.1,0) -- (10.1,0) node[right=2pt, font=\small\bfseries] {$H_C = 0$};
            \draw[-{Latex[length=2.5mm, width=1.8mm]}, line width=1.5pt, green!60!black] ($(9.1,0.7)+(-40:0.5)$) arc (-40:220:0.5);
            \node[above=6pt, green!60!black, font=\small\bfseries] at (9.1,0.7) {$M_C = 161.5\text{ kN}\cdot\text{m}$ (horario)};
        \end{tikzpicture}%
    }
    \caption{Diagrama de Cuerpo Libre (DCL) con resultantes y reacciones calculadas.}
\end{figure}

\paso{3}{Método de los cortes: leyes analíticas por tramos}

La viga queda dividida en 4 tramos por la presencia de cargas puntuales, rótulas y transiciones de carga:
\begin{enumerate}
    \item \textbf{Tramo 1:} $x \in [0, 1\text{ m}]$ (desde el apoyo $A$ hasta la carga puntual de $5\text{ kN}$).
    \item \textbf{Tramo 2:} $x \in [1, 2\text{ m}]$ (desde la carga de $5\text{ kN}$ hasta la rótula $B$).
    \item \textbf{Tramo 3:} $x \in [2, 3\text{ m}]$ (desde la rótula $B$ hasta el inicio de la carga triangular).
    \item \textbf{Tramo 4:} $x \in [3, 9\text{ m}]$ (tramo con carga triangular decreciente hasta el empotramiento $C$).
\end{enumerate}

\tramobadge{Tramo 1}{$x \in [0, 1\text{ m}]$}
Realizamos un corte a una distancia genérica $x \in [0, 1\text{ m}]$ reteniendo la porción izquierda:
\begin{align*}
    (\rightarrow +) \quad &\sum F_x = 0 \implies \boxed{N(x) = 0} \\
    (\uparrow +) \quad &\sum F_y = 0 \implies V_A - Q(x) = 0 \implies \boxed{Q(x) = 2.50\text{ kN}} \\
    (\curvearrowleft +) \quad &\sum M_{\text{corte}} = 0 \implies -V_A \cdot x + M_f(x) = 0 \implies \boxed{M_f(x) = 2.50\,x\text{ kN}\cdot\text{m}}
\end{align*}
\textbf{Valores en los extremos del tramo:}
\begin{itemize}
    \item En $x = 0$: $Q(0) = 2.50\text{ kN}$, \quad $M_f(0) = 0$.
    \item En $x = 1^-\text{ m}$: $Q(1^-) = 2.50\text{ kN}$, \quad $M_f(1) = 2.50\text{ kN}\cdot\text{m}$.
\end{itemize}

\tramobadge{Tramo 2}{$x \in [1, 2\text{ m}]$}
Cortamos reteniendo la porción izquierda, incorporando la carga puntual descendente en $x=1\text{ m}$:
\begin{align*}
    (\rightarrow +) \quad &\sum F_x = 0 \implies \boxed{N(x) = 0} \\
    (\uparrow +) \quad &\sum F_y = 0 \implies V_A - 5 - Q(x) = 0 \implies Q(x) = 2.50 - 5 \implies \boxed{Q(x) = -2.50\text{ kN}} \\
    (\curvearrowleft +) \quad &\sum M_{\text{corte}} = 0 \implies -V_A \cdot x + 5 \cdot (x - 1) + M_f(x) = 0 \\
    &\implies M_f(x) = 2.50\,x - 5(x - 1) = -2.50\,x + 5.00\text{ kN}\cdot\text{m}
\end{align*}
\textbf{Valores en los extremos del tramo:}
\begin{itemize}
    \item En $x = 1^+\text{ m}$: $Q(1^+) = -2.50\text{ kN}$ (salto $\Delta Q = -5\text{ kN}$ debido a la carga puntual), \quad $M_f(1) = 2.50\text{ kN}\cdot\text{m}$.
    \item En $x = 2^-\text{ m}$ (rótula $B$): $Q(2^-) = -2.50\text{ kN}$, \quad $M_f(2) = -2.50(2) + 5 = 0\text{ kN}\cdot\text{m}$ (¡se anula idénticamente en la rótula!).
\end{itemize}

\tramobadge{Tramo 3}{$x \in [2, 3\text{ m}]$}
En el tramo comprendido entre la rótula y el inicio de la carga triangular no hay nuevas cargas exteriores. Por tanto, los esfuerzos se mantienen idénticos a los del tramo 2:
\begin{align*}
    \boxed{N(x) = 0}, \qquad \boxed{Q(x) = -2.50\text{ kN}}, \qquad \boxed{M_f(x) = -2.50\,x + 5.00\text{ kN}\cdot\text{m}}
\end{align*}
\textbf{Valores en los extremos del tramo:}
\begin{itemize}
    \item En $x = 2^+\text{ m}$: $Q(2) = -2.50\text{ kN}$, \quad $M_f(2) = 0\text{ kN}\cdot\text{m}$.
    \item En $x = 3^-\text{ m}$: $Q(3^-) = -2.50\text{ kN}$, \quad $M_f(3) = -2.50(3) + 5 = -2.50\text{ kN}\cdot\text{m}$.
\end{itemize}

\tramobadge{Tramo 4}{$x \in [3, 9\text{ m}]$}
Para analizar este tramo con carga triangular que decrece hacia el empotramiento, resulta especialmente claro y ventajoso cortar la viga en $x$ y \textbf{retener la porción derecha} (entre la sección $x$ y el empotramiento $C$ en $x=9\text{ m}$), cuya longitud es $u = 9 - x$.

\textbf{1. Ley de carga triangular por semejanza geométrica:}
El triángulo se anula en $x=9$ y alcanza $q = 12\text{ kN/m}$ en $x=3$. A una distancia $u = 9 - x$ medida desde el extremo derecho $C$:
\[
    \frac{q(u)}{u} = \frac{12}{6} = 2 \implies q(u) = 2\,u = 2(9 - x)
\]
\textbf{2. Resultante y baricentro de la porción retenida derecha:}
\begin{itemize}
    \item Carga resultante del triángulo de base $u$:
    \[
        R(u) = \frac{1}{2} \cdot q(u) \cdot u = \frac{1}{2} \cdot 2u \cdot u = u^2 = (9 - x)^2
    \]
    \item Su punto de aplicación dista $\frac{u}{3} = \frac{9-x}{3}$ respecto al corte $x$ (ya que la máxima altura de este sub-triángulo está en la sección del corte $x$).
\end{itemize}

\textbf{3. Ecuaciones de equilibrio de la porción derecha:}
En la cara izquierda de la porción derecha retenida, el convenio de signos establece: $N$ hacia la izquierda, $Q$ hacia arriba ($\uparrow$) y $M_f$ en sentido horario ($\curvearrowright$). En el empotramiento $C$ actúan $V_C = 38.50\text{ kN}$ hacia arriba y el par de empotramiento $M_C = 161.50\text{ kN}\cdot\text{m}$ en sentido horario.

\begin{itemize}
    \item \textbf{Esfuerzo cortante $Q(x)$:}
    \begin{align}
        (\uparrow +) \quad \sum F_y = 0 &\implies Q(x) - R(u) + V_C = 0 \nonumber \\
        &\implies Q(x) - (9 - x)^2 + 38.50 = 0 \nonumber \\
        &\implies \boxed{Q(x) = (9 - x)^2 - 38.50 = x^2 - 18x + 42.50\text{ kN}}
    \end{align}
    \item \textbf{Momento flector $M_f(x)$:}
    Tomamos momentos respecto al punto del corte $x$ ($\curvearrowleft +$):
    \begin{align}
        \sum M_{\text{corte}} = 0 &\implies -M_f(x) - R(u) \cdot \frac{u}{3} + V_C \cdot u - M_C = 0 \nonumber \\
        &\implies M_f(x) = -\frac{u^3}{3} + 38.50\,u - 161.50 \nonumber \\
        &\implies \boxed{M_f(x) = -\frac{(9 - x)^3}{3} + 38.50(9 - x) - 161.50\text{ kN}\cdot\text{m}}
    \end{align}
\end{itemize}

\textbf{Valores notables y comprobación en los extremos del tramo 4:}
\begin{itemize}
    \item En $x = 3\text{ m}$ ($u = 6\text{ m}$):
    \begin{align*}
        Q(3) &= 6^2 - 38.50 = 36 - 38.50 = -2.50\text{ kN} \quad (\text{¡coincide con el valor al final del tramo 3!}) \\
        M_f(3) &= -\frac{6^3}{3} + 38.50(6) - 161.50 = -72 + 231 - 161.50 = -2.50\text{ kN}\cdot\text{m} \quad (\text{¡continuo!})
    \end{align*}
    \item En $x = 9\text{ m}$ ($u = 0$, empotramiento $C$):
    \begin{align*}
        Q(9) &= 0^2 - 38.50 = -38.50\text{ kN} \quad (\text{se equilibra con la reacción } V_C = +38.50\text{ kN}) \\
        M_f(9) &= 0 + 0 - 161.50 = -161.50\text{ kN}\cdot\text{m} \quad (\text{cierra exactamente en el valor de } -M_C)
    \end{align*}
\end{itemize}

\paso{4}{Diagramas de esfuerzos internos ($N, Q, M_f$)}

\begin{figure}[H]
    \centering
    \adjustbox{max width=0.98\textwidth}{%
        \begin{tikzpicture}[scale=1.2, >=Latex]
            % Líneas de referencia verticales
            \draw[dotted, line width=0.7pt, gray] (0, 0.5) -- (0, -8.2);
            \draw[dotted, line width=0.7pt, gray] (1, 0.5) -- (1, -8.2);
            \draw[dotted, line width=0.7pt, gray] (2, 0.5) -- (2, -8.2);
            \draw[dotted, line width=0.7pt, gray] (3, 0.5) -- (3, -8.2);
            \draw[dotted, line width=0.7pt, gray] (9, 0.5) -- (9, -8.2);

            % --- DIAGRAMA N(x) ---
            \node[left=8pt, font=\small\bfseries] at (0,0) {$N$};
            \draw[line width=1.2pt] (0,0) -- (9,0);
            \node[above=2pt, font=\footnotesize] at (4.5,0) {$N = 0$};

            % --- DIAGRAMA Q(x) ---
            \draw[line width=0.8pt, gray!60] (0,-2.6) -- (9,-2.6);
            \node[left=8pt, font=\small\bfseries] at (0,-2.6) {$Q$};

            % Tramo 1: Q = +2.5 kN
            \fill[blue!15] (0,-2.6) -- (0,-2.6+0.5) -- (1,-2.6+0.5) -- (1,-2.6) -- cycle;
            \draw[line width=1.3pt, blue!70!black] (0,-2.6+0.5) -- (1,-2.6+0.5);
            \draw[line width=1.1pt, blue!70!black] (0,-2.6) -- (0,-2.6+0.5);
            \node[above=1pt, blue!90!black, font=\footnotesize\bfseries] at (0.5,-2.6+0.5) {$+2.5$};

            % Salto en x=1m de +2.5 a -2.5
            \draw[line width=1.1pt, blue!70!black] (1,-2.6+0.5) -- (1,-2.6-0.5);

            % Tramo 2 y 3: Q = -2.5 kN constante de x=1 a x=3
            \fill[blue!15] (1,-2.6) -- (1,-2.6-0.5) -- (3,-2.6-0.5) -- (3,-2.6) -- cycle;
            \draw[line width=1.3pt, blue!70!black] (1,-2.6-0.5) -- (3,-2.6-0.5);
            \node[below=1pt, blue!90!black, font=\footnotesize\bfseries] at (2,-2.6-0.5) {$-2.5$};

            % Tramo 4: curva parabólica de Q de -2.5 a -38.5 en x=9
            \fill[blue!15] (3,-2.6) -- (3,-2.6-0.5) to[out=-25, in=170] (9,-2.6-1.8) -- (9,-2.6) -- cycle;
            \draw[line width=1.3pt, blue!70!black] (3,-2.6-0.5) to[out=-25, in=170] (9,-2.6-1.8);
            \draw[line width=1.1pt, blue!70!black] (9,-2.6-1.8) -- (9,-2.6);
            \node[right=3pt, blue!90!black, font=\footnotesize\bfseries] at (9,-2.6-1.8) {$-38.50\text{ kN}$};

            % Criterio de signos cortante negativo en vano 4
            \draw[-{Latex[length=2mm, width=1.4mm]}, line width=0.8pt, blue!80!black] (6.0,-3.2) -- (6.0,-3.7);
            \draw[-{Latex[length=2mm, width=1.4mm]}, line width=0.8pt, blue!80!black] (6.4,-3.7) -- (6.4,-3.2);

            % --- DIAGRAMA M_f(x) --- (Convenio: M_f negativo hacia arriba)
            \draw[line width=0.8pt, gray!60] (0,-6.0) -- (9,-6.0);
            \node[left=8pt, font=\small\bfseries] at (0,-6.0) {$M_f$};

            % Tramo 1: lineal de 0 a +2.5 en x=1 (positivo dibujado abajo)
            \fill[green!15] (0,-6.0) -- (1,-6.0-0.5) -- (1,-6.0) -- cycle;
            \draw[line width=1.3pt, green!50!black] (0,-6.0) -- (1,-6.0-0.5);
            \node[below=1pt, green!60!black, font=\footnotesize\bfseries] at (1,-6.0-0.5) {$+2.5$};

            % Tramo 2: lineal de +2.5 a 0 en x=2 (rótula)
            \fill[green!15] (1,-6.0) -- (1,-6.0-0.5) -- (2,-6.0) -- cycle;
            \draw[line width=1.3pt, green!50!black] (1,-6.0-0.5) -- (2,-6.0);
            \node[above=2pt, font=\scriptsize\bfseries] at (2,-6.0) {$0$ (Rótula)};

            % Tramo 3: lineal de 0 a -2.5 en x=3 (negativo dibujado arriba)
            \fill[blue!15] (2,-6.0) -- (3,-6.0+0.5) -- (3,-6.0) -- cycle;
            \draw[line width=1.3pt, blue!70!black] (2,-6.0) -- (3,-6.0+0.5);
            \node[above=1pt, blue!90!black, font=\footnotesize\bfseries] at (3,-6.0+0.5) {$-2.5$};

            % Tramo 4: cúbica de -2.5 en x=3 a -161.5 en x=9 (dibujada hacia arriba)
            \fill[blue!15] (3,-6.0) -- (3,-6.0+0.5) to[out=25, in=-165] (9,-6.0+2.2) -- (9,-6.0) -- cycle;
            \draw[line width=1.3pt, blue!70!black] (3,-6.0+0.5) to[out=25, in=-165] (9,-6.0+2.2);
            \draw[line width=1.1pt, blue!70!black] (9,-6.0+2.2) -- (9,-6.0);
            \node[right=3pt, blue!90!black, font=\footnotesize\bfseries] at (9,-6.0+2.2) {$-161.50\text{ kN}\cdot\text{m}$};

            % Criterio signo momento negativo arriba
            \draw[-{Stealth[length=2mm, width=1.4mm]}, line width=0.8pt, blue!80!black] ($(6.2,-5.3)+(90:0.25)$) arc (90:270:0.25);
            \node[font=\small\bfseries, blue!80!black] at (6.55,-5.3) {$-$};
            \draw[-{Stealth[length=2mm, width=1.4mm]}, line width=0.8pt, blue!80!black] ($(6.9,-5.3)+(90:0.25)$) arc (90:-90:0.25);
        \end{tikzpicture}%
    }
    \caption{Diagramas de esfuerzo axil ($N$), esfuerzo cortante ($Q$) y momento flector ($M_f$) con convenio de tracciones.}
    \label{fig:ejemplo21_diagramas}
\end{figure}

\comp{
    \textbf{Comprobación sistemática de los diagramas y equilibrio interno:}
    \begin{enumerate}
        \item \textbf{Relaciones diferenciales:}
        \begin{itemize}
            \item En tramos sin carga repartida ($x \in [0, 3\text{ m}]$): $\frac{dQ}{dx} = 0 \implies Q = \text{cte}$, y $\frac{dM_f}{dx} = Q$ (pendientes rectas coincidentes con el valor del cortante).
            \item En el tramo 4 ($x \in [3, 9\text{ m}]$): $\frac{dQ}{dx} = 2x - 18 = -2(9 - x) = -q(x)$, cumpliendo rigurosamente $\frac{dQ}{dx} = -q(x)$. Asimismo, $\frac{dM_f}{dx} = Q(x)$ y $\frac{d^2M_f}{dx^2} = -q(x) < 0$, confirmando concavidad numérica negativa.
        \end{itemize}
        \item \textbf{Condición de rótula interna ($x = 2\text{ m}$):}
        El momento flector evaluado en la rótula se anula estrictamente: $M_f(2) = 0\text{ kN}\cdot\text{m}$.
        \item \textbf{Condiciones de contorno y saltos:}
        \begin{itemize}
            \item Apoyo simple en $A$ ($x=0$): $Q(0^+) = V_A = 2.50\text{ kN}$, $M_f(0) = 0$.
            \item Carga puntual en $x = 1\text{ m}$: Discontinuidad de cortante $\Delta Q = Q(1^+) - Q(1^-) = -2.50 - 2.50 = -5.00\text{ kN} = -F$. El momento flector es continuo pero presenta un punto anguloso en el máximo positivo local $M_f(1) = +2.50\text{ kN}\cdot\text{m}$.
            \item Empotramiento en $C$ ($x = 9\text{ m}$): $Q(9^-) = -38.50\text{ kN}$, que cierra a cero con la reacción vertical $V_C = +38.50\text{ kN}$. El momento flector $M_f(9) = -161.50\text{ kN}\cdot\text{m}$ se equilibra exactamente con el momento reactivo horario $M_C = 161.50\text{ kN}\cdot\text{m}$.
        \end{itemize}
    \end{enumerate}
}

\nt{
    \textbf{Detalle conceptual del empotramiento:}
    En una sección de empotramiento rígido, la rotación de la directriz está totalmente coaccionada ($\theta = 0$). El momento flector negativo obtenido ($-161.50\text{ kN}\cdot\text{m}$) indica que las fibras superiores del perfil resultan sometidas a tracción, requiriendo un par reactivo en sentido horario para equilibrar las cargas actuantes del vano.
}

\subsection{Ejemplo 22: Viga empotrada con carga triangular, rótula interna, apoyo móvil intermedio y voladizo con carga puntual}

\begin{tcolorbox}[colback=blue!3!white, colframe=blue!40!black, arc=3pt, boxrule=0.8pt, left=8pt, right=8pt, top=6pt, bottom=6pt]
    \centering
    {\Large\bfseries\color{blue!60!black} Viga con empotramiento, rótula interna, apoyo intermedio y voladizo}\\[2pt]
    {\normalsize\bfseries Viga de $L = 7\text{ m}$ empotrada en $A$ ($x=0$), rótula en $x=4\text{ m}$, apoyo móvil en $B$ ($x=5\text{ m}$) y carga puntual en el extremo libre $C$ ($x=7\text{ m}$)}

    \vspace{4pt}\hrule\vspace{4pt}
    {\small\textbf{Unidades de trabajo:} $L\,[\text{m}] \quad F\,[\text{kN}] \quad M\,[\text{kN}\cdot\text{m}] \quad q\,[\text{kN/m}]$}
\end{tcolorbox}

\textbf{Enunciado:}
Sea una viga recta de longitud total $L = 7\text{ m}$ empotrada en su extremo izquierdo $A$ ($x = 0$), con una rótula o articulación interna en la sección $x = 4\text{ m}$, sustentada por un apoyo móvil en $B$ ($x = 5\text{ m}$) y con un voladizo que termina en el extremo libre $C$ ($x = 7\text{ m}$).

Sobre la estructura actúan las siguientes solicitaciones exteriores:
\begin{itemize}
    \item Una carga distribuida triangular descendente que crece linealmente desde $0$ en el empotramiento $A$ ($x=0$) hasta un valor máximo $q_{\text{máx}} = 2\text{ kN/m}$ en $x = 3\text{ m}$ (base de $3\text{ m}$).
    \item Un tramo sin carga repartida entre $x = 3\text{ m}$ y $x = 7\text{ m}$.
    \item Una carga puntual vertical descendente $F = 10\text{ kN}$ aplicada en el extremo libre derecho $C$ ($x = 7\text{ m}$).
\end{itemize}

Se pide:
\begin{enumerate}[label=(\alph*)]
    \item Clasificación estática del sistema.
    \item Diagrama de Cuerpo Libre (DCL) y cálculo de las reacciones en los apoyos.
    \item Determinación de las leyes analíticas de esfuerzos internos ($N(x), Q(x), M_f(x)$) por el método de cortes.
    \item Trazado de los diagramas de solicitaciones internas acotando valores singulares y verificando el convenio de signos.
    \item Comprobaciones sistemáticas de equilibrio y contorno.
\end{enumerate}

\begin{figure}[H]
    \centering
    \adjustbox{max width=0.98\textwidth}{%
        \begin{tikzpicture}[scale=1.3, >=Latex]
            % Viga principal
            \draw[line width=2.4pt, blue!80!black] (0,0) -- (7,0);

            % Cotas inferiores: 3m, 1m, 1m, 2m
            \draw[|<->|, >=Latex] (0,-1.1) -- (3,-1.1) node[midway, fill=white, font=\footnotesize] {$3\text{ m}$};
            \draw[|<->|, >=Latex] (3,-1.1) -- (4,-1.1) node[midway, fill=white, font=\footnotesize] {$1\text{ m}$};
            \draw[|<->|, >=Latex] (4,-1.1) -- (5,-1.1) node[midway, fill=white, font=\footnotesize] {$1\text{ m}$};
            \draw[|<->|, >=Latex] (5,-1.1) -- (7,-1.1) node[midway, fill=white, font=\footnotesize] {$2\text{ m}$};

            % Empotramiento en A (x = 0)
            \coordinate (A) at (0,0);
            \draw[line width=1.8pt] (0,-0.8) -- (0,1.2);
            \foreach \y in {-0.7,-0.4,-0.1,0.2,0.5,0.8,1.1} {
                \draw[thick] (0,\y) -- (-0.28,\y+0.2);
            }
            \node[above left=2pt, font=\bfseries] at (A) {$A$};

            % Carga triangular 0 -> 2 kN/m en [0, 3m] (h = 1.3)
            \draw[line width=1.1pt, red!80!black] (0,0) -- (3,1.3) -- (3,0);
            \foreach \x in {0.6, 1.2, 1.8, 2.4, 3.0} {
                \pgfmathsetmacro{\h}{\x*(1.3/3.0)}
                \draw[-{Latex[length=2.2mm, width=1.4mm]}, line width=0.8pt, red!80!black] (\x, \h) -- (\x, 0.05);
            }
            \node[above=2pt, red!80!black, font=\small\bfseries] at (3.0,1.3) {$2\text{ kN/m}$};

            % Rótula interna en x = 4m
            \filldraw[fill=white, draw=black, line width=1.4pt] (4,0) circle (3pt) node[above=5pt, font=\footnotesize\bfseries] {Rótula ($x=4\text{ m}$)};

            % Apoyo móvil en B (x = 5m)
            \coordinate (B) at (5,0);
            \fill (B) circle (2pt) node[above right=2pt, font=\bfseries] {$B$};
            \draw[thick] (B) -- ++(-0.28,-0.4) -- ++(0.56,0) -- cycle;
            \draw[thick] ($(B)+(-0.16,-0.5)$) circle (0.09);
            \draw[thick] ($(B)+(0.16,-0.5)$) circle (0.09);
            \draw[line width=1pt] ($(B)+(-0.35,-0.61)$) -- ($(B)+(0.35,-0.61)$);

            % Carga puntual en C (x = 7m)
            \draw[-{Latex[length=3.2mm, width=2.2mm]}, line width=1.5pt, red!85!black] (7,1.4) -- (7,0.05);
            \node[above, red!85!black, font=\small\bfseries] at (7,1.4) {$10\text{ kN}$};
            \fill (7,0) circle (2pt) node[below right=2pt, font=\bfseries] {$C$};
        \end{tikzpicture}%
    }
    \caption{Esquema geométrico y de cargas de la viga.}
    \label{fig:ejemplo22_esquema}
\end{figure}

\paso{1}{Datos y clasificación estática}

Determinamos las condiciones de sustentación y vinculación del sistema:
\begin{itemize}
    \item Número de barras continuas: $m = 1$.
    \item Número de reacciones incógnita: $r = 4$ ($H_A, V_A, M_A$ en el empotramiento izquierdo $A$, y $V_B$ en el apoyo móvil intermedio).
    \item Grados de libertad liberados por enlaces internos: $h = 1$ (la rótula en $x = 4\text{ m}$ permite la libre rotación relativa, por lo que $M_f(4) = 0$).
    \item Grado de hiperestaticidad:
    \[
        GH = r - 3 - h = 4 - 3 - 1 = 0 \implies \text{\textbf{Estructura isostática.}}
    \]
\end{itemize}

\paso{2}{Diagrama de Cuerpo Libre (DCL) y Cálculo de Reacciones}

\paragraph{1. Aislamiento de la subestructura derecha ($x \in [4, 7\text{ m}]$):}
La porción derecha comprendida entre la rótula ($x = 4\text{ m}$) y el extremo libre ($x = 7\text{ m}$) constituye una viga simplemente apoyada en $B$ ($x = 5\text{ m}$) y apoyada en la rótula en $x = 4\text{ m}$. En la rótula actúa la fuerza vertical interna $V_{\text{rót}}$ transmitida, y el momento flector se anula ($M_f(4) = 0$).

Tomando momentos respecto a la rótula ($x = 4\text{ m}$) para la porción derecha ($\curvearrowleft +$):
\begin{equation}
    \sum M_{\text{rótula}}^{\text{der}} = 0 \implies V_B \cdot (5 - 4) - 10 \cdot (7 - 4) = 0 \implies V_B \cdot 1 - 30 = 0 \implies \boxed{V_B = 30.00\text{ kN}}
\end{equation}
Por equilibrio vertical de esta subestructura derecha:
\begin{equation}
    \sum F_{y,\text{der}} = 0 \implies V_{\text{rót}} + V_B - 10 = 0 \implies V_{\text{rót}} + 30 - 10 = 0 \implies V_{\text{rót}} = -20\text{ kN}
\end{equation}
Por tanto, la porción derecha ejerce sobre la porción izquierda en la rótula una fuerza vertical descendente de $20\text{ kN}$ (esfuerzo cortante $Q(4) = -20\text{ kN}$).

\paragraph{2. Resultante de la carga triangular ($x \in [0, 3\text{ m}]$):}
\[
    R_{\text{tri}} = \frac{1}{2} \cdot q_{\text{máx}} \cdot b = \frac{1}{2} \cdot 2 \cdot 3 = 3\text{ kN}
\]
Su punto de aplicación baricéntrico se sitúa a $\frac{2}{3}$ de la base desde el vértice $A$ ($x = 0$):
\[
    x_G = \frac{2}{3} \cdot 3 = 2\text{ m desde } A \quad (1\text{ m a la izquierda de } x=3)
\]

\paragraph{3. Equilibrio global de la viga completa:}
\begin{itemize}
    \item \textbf{Equilibrio horizontal:}
    \begin{equation}
        (\rightarrow +) \quad \sum F_x = 0 \implies \boxed{H_A = 0}
    \end{equation}
    \item \textbf{Equilibrio vertical:}
    \begin{equation}
        (\uparrow +) \quad \sum F_y = 0 \implies V_A - R_{\text{tri}} + V_B - 10 = 0 \implies V_A - 3 + 30 - 10 = 0 \implies V_A + 17 = 0 \implies \boxed{V_A = -17.00\text{ kN}}
    \end{equation}
    El signo negativo indica que la reacción vertical en el empotramiento $A$ actúa hacia abajo ($\downarrow$), impidiendo el levantamiento de la viga.
    \item \textbf{Equilibrio de momentos respecto al empotramiento $A$ ($x = 0$, $\curvearrowleft +$):}
    \begin{align}
        \sum M_A = 0 &\implies M_A - R_{\text{tri}} \cdot 2 + V_B \cdot 5 - 10 \cdot 7 = 0 \nonumber \\
        &\implies M_A - 3 \cdot 2 + 30 \cdot 5 - 70 = 0 \nonumber \\
        &\implies M_A - 6 + 150 - 70 = 0 \implies M_A + 74 = 0 \implies \boxed{M_A = -74.00\text{ kN}\cdot\text{m}}
    \end{align}
    Por consiguiente, el momento de empotramiento reactivo actúa en \textbf{sentido horario} con módulo $74\text{ kN}\cdot\text{m}$ (traccionando las fibras superiores de la viga).
\end{itemize}

\begin{figure}[H]
    \centering
    \adjustbox{max width=0.98\textwidth}{%
        \begin{tikzpicture}[scale=1.2, >=Latex]
            % Viga DCL
            \draw[line width=2.4pt, blue!80!black] (0,0) -- (7,0);

            % Reacción vertical V_A hacia abajo
            \draw[->, >=Latex, line width=1.5pt, green!60!black] (0,0) -- (0,-1.3) node[below=2pt, font=\small\bfseries] {$V_A = 17\text{ kN} \ (\downarrow)$};
            % Momento horario de empotramiento en A
            \draw[-{Latex[length=2.5mm, width=1.8mm]}, line width=1.5pt, green!60!black] ($(0,0.6)+(140:0.45)$) arc (140:-140:0.45);
            \node[above=4pt, green!60!black, font=\small\bfseries] at (0,0.6) {$M_A = 74\text{ kN}\cdot\text{m}$ (horario)};

            % Resultante triangular
            \draw[-{Latex[length=2.8mm, width=2mm]}, line width=1.4pt, orange!90!black] (2,1.3) -- (2,0.05);
            \node[above, orange!90!black, font=\small\bfseries] at (2,1.3) {$R_{\text{tri}} = 3\text{ kN}$};

            % Rótula
            \filldraw[fill=white, draw=black, line width=1.3pt] (4,0) circle (3pt) node[above=3pt, font=\footnotesize\bfseries] {Rótula};

            % Reacción V_B hacia arriba
            \draw[->, >=Latex, line width=1.5pt, green!60!black] (5,-1.3) -- (5,-0.05) node[midway, right=2pt, font=\small\bfseries] {$V_B = 30\text{ kN} \ (\uparrow)$};

            % Carga puntual en C
            \draw[-{Latex[length=2.8mm, width=2mm]}, line width=1.4pt, red!85!black] (7,1.3) -- (7,0.05);
            \node[above, red!85!black, font=\small\bfseries] at (7,1.3) {$10\text{ kN}$};

            % Cotas
            \draw[|<->|, >=Latex, gray] (0,-0.6) -- (2,-0.6) node[midway, fill=white, font=\tiny] {$2\text{ m}$};
            \draw[|<->|, >=Latex, gray] (2,-0.6) -- (3,-0.6) node[midway, fill=white, font=\tiny] {$1\text{ m}$};
            \draw[|<->|, >=Latex, gray] (3,-0.6) -- (4,-0.6) node[midway, fill=white, font=\tiny] {$1\text{ m}$};
            \draw[|<->|, >=Latex, gray] (4,-0.6) -- (5,-0.6) node[midway, fill=white, font=\tiny] {$1\text{ m}$};
            \draw[|<->|, >=Latex, gray] (5,-0.6) -- (7,-0.6) node[midway, fill=white, font=\tiny] {$2\text{ m}$};
        \end{tikzpicture}%
    }
    \caption{Diagrama de Cuerpo Libre (DCL) con reacciones y resultantes.}
    \label{fig:ejemplo22_dcl}
\end{figure}

\paso{3}{Método de los cortes: leyes analíticas por tramos}

La estructura presenta 4 tramos delimitados por el fin de la carga triangular ($x=3\text{ m}$), la rótula interna ($x=4\text{ m}$) y el apoyo intermedio ($x=5\text{ m}$):

\tramobadge{Tramo 1}{$x \in [0, 3\text{ m}]$}
Porción izquierda con la sección de empotramiento $A$ ($V_A = 17\text{ kN}\downarrow$, par horario $M_A = 74\text{ kN}\cdot\text{m}$).
La carga triangular crece linealmente con pendiente $m = \frac{2}{3}\text{ kN/m}^2$:
\[
    q(x) = \frac{2}{3}\,x, \qquad R(x) = \frac{1}{2} \cdot x \cdot \left(\frac{2}{3}x\right) = \frac{x^2}{3}, \quad \text{a distancia } \frac{x}{3} \text{ del corte.}
\]
Planteando el equilibrio en el corte:
\begin{align*}
    (\rightarrow +) \quad &\sum F_x = 0 \implies \boxed{N(x) = 0} \\
    (\uparrow +) \quad &\sum F_y = 0 \implies -V_A - R(x) - Q(x) = 0 \implies -17 - \frac{x^2}{3} - Q(x) = 0 \\
    &\implies \boxed{Q(x) = -17 - \frac{x^2}{3}\text{ kN}} \\
    (\curvearrowleft +) \quad &\sum M_{\text{corte}} = 0 \implies -M_A + V_A \cdot x + R(x) \cdot \frac{x}{3} + M_f(x) = 0 \\
    &\implies -74 + 17\,x + \left(\frac{x^2}{3}\right)\left(\frac{x}{3}\right) + M_f(x) = 0 \\
    &\implies \boxed{M_f(x) = 74 - 17\,x - \frac{x^3}{9}\text{ kN}\cdot\text{m}}
\end{align*}
\textbf{Valores en los extremos del tramo:}
\begin{itemize}
    \item En $x = 0$: $Q(0) = -17.00\text{ kN}$, \quad $M_f(0) = 74.00\text{ kN}\cdot\text{m}$ (según convenio gráfico, tracción superior $\implies$ dibujado hacia arriba en $74$).
    \item En $x = 3\text{ m}$: $Q(3) = -17 - \frac{9}{3} = -20.00\text{ kN}$, \quad $M_f(3) = 74 - 17(3) - \frac{27}{9} = 74 - 51 - 3 = 20.00\text{ kN}\cdot\text{m}$.
\end{itemize}

\tramobadge{Tramo 2}{$x \in [3, 4\text{ m}]$}
Entre el fin de la carga triangular y la rótula interna no hay cargas intermedias ($q = 0$). La porción izquierda contiene la resultante total $R_{\text{tri}} = 3\text{ kN}$ en $x = 2\text{ m}$:
\begin{align*}
    (\uparrow +) \quad &\sum F_y = 0 \implies -17 - 3 - Q(x) = 0 \implies \boxed{Q(x) = -20.00\text{ kN}} \ (\text{constante}) \\
    (\curvearrowleft +) \quad &\sum M_{\text{corte}} = 0 \implies -74 + 17\,x + 3\,(x - 2) + M_f(x) = 0 \\
    &\implies -74 - 6 + 20\,x + M_f(x) = 0 \implies \boxed{M_f(x) = 80 - 20\,x\text{ kN}\cdot\text{m}}
\end{align*}
\textbf{Valores en los extremos del tramo:}
\begin{itemize}
    \item En $x = 3\text{ m}$: $Q(3) = -20.00\text{ kN}$, \quad $M_f(3) = 80 - 20(3) = 20.00\text{ kN}\cdot\text{m}$ (continuo).
    \item En $x = 4\text{ m}$ (rótula): $Q(4^-) = -20.00\text{ kN}$, \quad $M_f(4) = 80 - 20(4) = 0\text{ kN}\cdot\text{m}$ (¡se anula exactamente en la rótula!).
\end{itemize}

\tramobadge{Tramo 3}{$x \in [4, 5\text{ m}]$}
En el tramo entre la rótula y el apoyo intermedio $B$, la ley de esfuerzos continúa idéntica al tramo anterior al no existir acciones aplicadas puntualmente en la rótula:
\begin{align*}
    \boxed{N(x) = 0}, \qquad \boxed{Q(x) = -20.00\text{ kN}}, \qquad \boxed{M_f(x) = 80 - 20\,x\text{ kN}\cdot\text{m}}
\end{align*}
\textbf{Valores en los extremos del tramo:}
\begin{itemize}
    \item En $x = 4^+\text{ m}$: $Q(4) = -20.00\text{ kN}$, \quad $M_f(4) = 0\text{ kN}\cdot\text{m}$.
    \item En $x = 5^-\text{ m}$ (justo a la izquierda del apoyo $B$):
    \[
        Q(5^-) = -20.00\text{ kN}, \qquad M_f(5) = 80 - 20(5) = -20.00\text{ kN}\cdot\text{m}
    \]
\end{itemize}

\tramobadge{Tramo 4}{$x \in [5, 7\text{ m}]$}
Para este tramo en voladizo resulta directo cortar y retener la \textbf{porción derecha} desde el corte hasta el extremo libre $C$ ($x = 7\text{ m}$), de longitud $7 - x$, donde solo actúa la carga puntual $F = 10\text{ kN}\downarrow$:
\begin{align*}
    (\uparrow +) \quad &Q(x) - 10 = 0 \implies \boxed{Q(x) = +10.00\text{ kN}} \ (\text{constante}) \\
    (\curvearrowleft +) \quad &-M_f(x) - 10 \cdot (7 - x) = 0 \implies \boxed{M_f(x) = -10\,(7 - x) = 10\,x - 70\text{ kN}\cdot\text{m}}
\end{align*}
\textbf{Valores en los extremos del tramo:}
\begin{itemize}
    \item En $x = 5^+\text{ m}$ (a la derecha de $B$):
    \[
        Q(5^+) = +10.00\text{ kN}, \qquad M_f(5) = -10(2) = -20.00\text{ kN}\cdot\text{m}
    \]
    El cortante experimenta una discontinuidad por la reacción de apoyo:
    \[
        \Delta Q = Q(5^+) - Q(5^-) = +10.00 - (-20.00) = +30.00\text{ kN} = V_B
    \]
    \item En $x = 7\text{ m}$ (extremo libre $C$):
    \[
        Q(7^-) = +10.00\text{ kN} \ (\text{cierra a cero con } F = 10\text{ kN}), \qquad M_f(7) = 0\text{ kN}\cdot\text{m}
    \]
\end{itemize}

\paso{4}{Diagramas de esfuerzos internos ($N, Q, M_f$)}

\begin{figure}[H]
    \centering
    \adjustbox{max width=0.98\textwidth}{%
        \begin{tikzpicture}[scale=1.2, >=Latex]
            % Líneas de referencia verticales alineadas
            \draw[dotted, line width=0.7pt, gray] (0, 0.5) -- (0, -8.2);
            \draw[dotted, line width=0.7pt, gray] (3, 0.5) -- (3, -8.2);
            \draw[dotted, line width=0.7pt, gray] (4, 0.5) -- (4, -8.2);
            \draw[dotted, line width=0.7pt, gray] (5, 0.5) -- (5, -8.2);
            \draw[dotted, line width=0.7pt, gray] (7, 0.5) -- (7, -8.2);

            % --- DIAGRAMA N(x) ---
            \node[left=8pt, font=\small\bfseries] at (0,0) {$N$};
            \draw[line width=1.2pt] (0,0) -- (7,0);
            \node[above=2pt, font=\footnotesize] at (3.5,0) {$N = 0$};

            % --- DIAGRAMA Q(x) ---
            \draw[line width=0.8pt, gray!60] (0,-2.6) -- (7,-2.6);
            \node[left=8pt, font=\small\bfseries] at (0,-2.6) {$Q$};

            % Tramo 1: parabólica de -17 a -20 en x=3
            \fill[blue!15] (0,-2.6) -- (0,-2.6-1.1) to[out=-15, in=170] (3,-2.6-1.35) -- (3,-2.6) -- cycle;
            \draw[line width=1.3pt, blue!70!black] (0,-2.6-1.1) to[out=-15, in=170] (3,-2.6-1.35);
            \draw[line width=1.1pt, blue!70!black] (0,-2.6) -- (0,-2.6-1.1);
            \node[left=2pt, blue!90!black, font=\footnotesize\bfseries] at (0,-2.6-1.1) {$-17$};
            \node[below=2pt, blue!90!black, font=\footnotesize\bfseries] at (3,-2.6-1.35) {$-20$};

            % Tramo 2 y 3: constante en -20 de x=3 a x=5
            \fill[blue!15] (3,-2.6) -- (3,-2.6-1.35) -- (5,-2.6-1.35) -- (5,-2.6) -- cycle;
            \draw[line width=1.3pt, blue!70!black] (3,-2.6-1.35) -- (5,-2.6-1.35);
            \node[below=2pt, blue!90!black, font=\footnotesize\bfseries] at (5,-2.6-1.35) {$-20$};

            % Criterio cortante negativo en tramo intermedio
            \draw[-{Latex[length=1.8mm, width=1.2mm]}, line width=0.75pt, blue!80!black] (3.8,-2.9) -- (3.8,-3.4);
            \draw[-{Latex[length=1.8mm, width=1.2mm]}, line width=0.75pt, blue!80!black] (4.2,-3.4) -- (4.2,-2.9);

            % Salto de cortante en B (x=5) de -20 a +10 (+30 kN)
            \draw[line width=1.3pt, blue!70!black] (5,-2.6-1.35) -- (5,-2.6+0.7);

            % Tramo 4: constante en +10 de x=5 a x=7
            \fill[blue!15] (5,-2.6) -- (5,-2.6+0.7) -- (7,-2.6+0.7) -- (7,-2.6) -- cycle;
            \draw[line width=1.3pt, blue!70!black] (5,-2.6+0.7) -- (7,-2.6+0.7);
            \draw[line width=1.1pt, blue!70!black] (7,-2.6+0.7) -- (7,-2.6);
            \node[above=2pt, blue!90!black, font=\footnotesize\bfseries] at (6,-2.6+0.7) {$+10$};
            % Criterio cortante positivo en voladizo
            \draw[-{Latex[length=1.8mm, width=1.2mm]}, line width=0.75pt, blue!80!black] (6.0,-2.4) -- (6.0,-1.9);
            \draw[-{Latex[length=1.8mm, width=1.2mm]}, line width=0.75pt, blue!80!black] (6.4,-1.9) -- (6.4,-2.4);

            % --- DIAGRAMA M_f(x) --- (Convenio: momentos negativos / tracción arriba graficados hacia arriba)
            \draw[line width=0.8pt, gray!60] (0,-6.0) -- (7,-6.0);
            \node[left=8pt, font=\small\bfseries] at (0,-6.0) {$M_f$};

            % Tramo 1: cúbica decreciente desde 74 hasta 20 en x=3 (dibujado arriba por tracción superior)
            \fill[orange!20] (0,-6.0) -- (0,-6.0+1.85) to[out=-25, in=170] (3,-6.0+0.5) -- (3,-6.0) -- cycle;
            \draw[line width=1.3pt, orange!85!black] (0,-6.0+1.85) to[out=-25, in=170] (3,-6.0+0.5);
            \draw[line width=1.1pt, orange!85!black] (0,-6.0) -- (0,-6.0+1.85);
            \node[left=2pt, orange!90!black, font=\footnotesize\bfseries] at (0,-6.0+1.85) {$74$};
            \node[above=2pt, orange!90!black, font=\footnotesize\bfseries] at (3,-6.0+0.5) {$20$};

            % Criterio momento tracción superior (arriba)
            \draw[-{Stealth[length=2mm, width=1.4mm]}, line width=0.8pt, orange!85!black] ($(1.6,-5.4)+(90:0.25)$) arc (90:270:0.25);
            \draw[-{Stealth[length=2mm, width=1.4mm]}, line width=0.8pt, orange!85!black] ($(2.2,-5.4)+(90:0.25)$) arc (90:-90:0.25);

            % Tramo 2: lineal desde 20 en x=3 hasta 0 en x=4 (rótula)
            \fill[orange!20] (3,-6.0) -- (3,-6.0+0.5) -- (4,-6.0) -- cycle;
            \draw[line width=1.3pt, orange!85!black] (3,-6.0+0.5) -- (4,-6.0);
            \node[above=3pt, font=\scriptsize\bfseries] at (4,-6.0) {$0$};

            % Tramo 3: lineal desde 0 en x=4 hasta -20 en x=5 (dibujado abajo / tracción inferior o cambia a vértice)
            \fill[orange!20] (4,-6.0) -- (5,-6.0-0.5) -- (5,-6.0) -- cycle;
            \draw[line width=1.3pt, orange!85!black] (4,-6.0) -- (5,-6.0-0.5);
            \node[below=2pt, orange!90!black, font=\footnotesize\bfseries] at (5,-6.0-0.5) {$-20$};

            % Tramo 4: lineal desde -20 en x=5 hasta 0 en x=7
            \fill[orange!20] (5,-6.0) -- (5,-6.0-0.5) -- (7,-6.0) -- cycle;
            \draw[line width=1.3pt, orange!85!black] (5,-6.0-0.5) -- (7,-6.0);
            \node[above=2pt, font=\scriptsize\bfseries] at (7,-6.0) {$0$};

            % Criterio momento en el vértice inferior sobre B
            \draw[-{Stealth[length=1.8mm, width=1.2mm]}, line width=0.75pt, orange!85!black] ($(5.0,-6.7)+(270:0.22)$) arc (270:90:0.22);
            \draw[-{Stealth[length=1.8mm, width=1.2mm]}, line width=0.75pt, orange!85!black] ($(5.6,-6.7)+(270:0.22)$) arc (270:450:0.22);
        \end{tikzpicture}%
    }
    \caption{Diagramas de esfuerzo axil ($N$), esfuerzo cortante ($Q$) y momento flector ($M_f$) acotando valores característicos.}
    \label{fig:ejemplo22_diagramas}
\end{figure}

\comp{
    \textbf{Comprobación sistemática de las leyes y diagramas:}
    \begin{enumerate}
        \item \textbf{Relaciones diferenciales:}
        \begin{itemize}
            \item Tramo 1 ($x \in [0, 3\text{ m}]$): $\frac{dQ}{dx} = -\frac{2}{3}x = -q(x)$, y $\frac{dM_f}{dx} = -17 - \frac{x^2}{3} = Q(x)$.
            \item Tramo 2 y 3 ($x \in [3, 5\text{ m}]$): $Q = -20\text{ kN} = \text{cte} \implies \frac{dM_f}{dx} = -20 = Q(x)$.
            \item Tramo 4 ($x \in [5, 7\text{ m}]$): $Q = +10\text{ kN} = \text{cte} \implies \frac{dM_f}{dx} = +10 = Q(x)$.
        \end{itemize}
        \item \textbf{Condición de rótula interna ($x = 4\text{ m}$):}
        El momento flector se anula estrictamente: $M_f(4) = 0\text{ kN}\cdot\text{m}$.
        \item \textbf{Condiciones de contorno y saltos:}
        \begin{itemize}
            \item Empotramiento en $A$ ($x = 0$): $Q(0) = -17.00\text{ kN} = V_A\downarrow$, y $M_f(0) = 74.00\text{ kN}\cdot\text{m}$ correspondiente con el par horario de empotramiento $M_A$.
            \item Apoyo móvil en $B$ ($x = 5\text{ m}$): Salto en cortante $\Delta Q = Q(5^+) - Q(5^-) = +10 - (-20) = 30\text{ kN} = V_B\uparrow$.
            \item Extremo libre $C$ ($x = 7\text{ m}$): $M_f(7) = 0$ y el cortante $Q(7) = +10\text{ kN}$ cierra con la carga puntual descendente $F = 10\text{ kN}$.
        \end{itemize}
    \end{enumerate}
}



